%!TEX root = ./../../springer-mono-format/book.tex
%!TEX root = ./../../thesis-trier.tex

Here, in \cref{constrsync:sec:codes:basic-notions} we introduce codes and overlap-free languages
and state some basic relations. Some of our later results will use prefix codes and solid codes. 
A few conditions using prefix codes use constructions related to maximal prefix codes. This
class of codes is discussed in \cref{constrsync:sec:max-prefix-codes}.
In \cref{constrsync:sec:more-on-prefix-codes} we collect properties of prefix codes that we need for those results.
\Cref{constrsync:sec:codes:levensthein-romanov} states some further characterizations of our classes of codes and overlap-free languages and relations to the Levenshtein-Romanov mapping.
\Cref{constrsync:sec:codes:solid} collects characterizations of solid codes (some of them have
not been observed before).

\subsection{Basic Notions}
\label{constrsync:sec:codes:basic-notions}

Here, in \cref{constrsync:def:codes}, we introduce overlap-free languages,
codes, and various subclasses of codes as, for example, solid codes, comma-free codes, and infix codes.
%Here, in \cref{constrsync:def:codes}, we introduce codes, various subclasses thereof as, for example, solid codes, comma-free codes, and infix codes,
%and overlap-free languages.
\Cref{constr_sync:prop:selfsync_implies_avoiding} 
proves that every solid code is a comma-free code, and that every comma-free code is an infix code.
In \cref{constrsync:ex:binarysolidcode} we also show that overlap-free languages
are not codes in general. The remaining inclusions among infix, bifix, prefix, and suffix codes
are obvious. \Cref{fig:code-hasse} gives an overview of the inclusion relations among
the introduces code classes and overlap-free languages.
As formally stated in \cref{constr_sync:prop:coding_morphisms},
codes are used to encode symbols from some source alphabet.
%Codes are used to encode symbol from some source alphabet,
%a formal statement of this fact is given
%in \cref{constr_sync:prop:coding_morphisms}.
%In \cref{constr_sync:prop:coding_morphisms} we formalize 
%the idea that codes encode symbols from a source alphabet.
%the relation
%between a source alphabet and codes as an encoding of symbols from such an alphabet.
\Cref{constr_sync:lem:two-words-code,constrsync:prop:conjugate-words} 
give two combinatorial results on word equations.

% L = emptyset erlauben? macht für präfixcode bei einigen aussagen sinn.
\begin{definition} \label{constrsync:def:codes}
	Let $L \subseteq \Sigma^+$ be non-empty and $u,v,w \in \Sigma^*$.
	
	\IN{code}
	\IN{overlap-free code}
	\IN{infix code}
	\IN{prefix code}
	\IN{suffix code}
	\IN{bifix code}
	\IN{solid code}
	\IN{comma-free code}
	\index{unique factorization property|see{code}}
	\begin{enumerate}%[itemsep=-2pt]	
		\item \label{constrsync:item:code} $L$ is a \emph{code} if, for $u_1, \ldots, u_n, v_1, \ldots, v_m \in L$,
		the equality
		\[
		u_1 \cdots u_n = v_1 \cdots v_m
		\]
		implies $n = m$ and $u_i = v_i$ for all $i \in \{1,\ldots,n\}$ (the \emph{unique factorization property} in $L^*$ w.r.t.~$L$).
		
		\item  $L$ is \emph{overlap-free}		
		if $uv, vw \in L$ with $u \ne\varepsilon$ and  $w \ne \varepsilon$ implies $v = \varepsilon$, i.e., no proper non-empty suffix
		of a word from~$L$ is also a proper non-empty prefix of some word in~$L$.\footnote{The condition
			that~$u$ and~$v$ are non-empty is missing in~\cite{DBLP:journals/tcs/JurgensenS21},
			but it is essential since without it,
			for $v \in L$, we have $\varepsilon\cdot v, v \cdot \varepsilon \in L$
			and thus $v = \varepsilon$, which
			contradicts the fact that $L$ is a subset of $\Sigma^+$.}		
		
		%\item $L$ is a \emph{block code} if it is a code and there exits $n > 0$ such that $L \subseteq \Sigma^n$.
		
		\item $L$ is an \emph{infix code} (or \emph{infix-free}) if $u, vuw \in L$ implies $v = w = \varepsilon$.
		
		\item $L$ is a \emph{prefix code} (or \emph{prefix-free}) if $uv, u \in L$ implies $v = \varepsilon$.
		
		\item $L$ is a \emph{suffix code} (or \emph{suffix-free}) if $vu, u \in L$ implies $v = \varepsilon$.
		
		\item $L$ is a \emph{bifix code} if it is a prefix and a suffix code. 
		
		\item $L$ is a \emph{solid code} if it is an overlap-free infix code.
		
		\item $L$ is a \emph{comma-free code} if $L^2 \cap \Sigma^+ L \Sigma^+ = \emptyset$.\footnote{The name comma-free refers to the fact that, when facing a stream of code words,
			we can uniquely decipher the next code words when 
			starting to read from somewhere in the middle of the stream
			without the need of a special separation symbol.}
	\end{enumerate}
\end{definition}

\Cref{constrsync:fig:freeness-fails} gives graphical depictions when some of the defining
conditions from \cref{constrsync:def:codes} fail.

\begin{table}[htbp]
	\centering
	\caption[Failure of defining conditions for various language classes]{Situations where the defining 
		conditions for overlap-, infix-, prefix-, suffix-, and comma-freeness fail. 
		Boxes labeled $u,v,w$ represent words or factors, and their relative positions indicate the overlaps.}
	\label{constrsync:fig:freeness-fails}
	
	\setlength{\aboverulesep}{6pt}
	\setlength{\belowrulesep}{6pt}
	
	\begin{tabular}{m{4cm} m{6cm}}
		\toprule
		not overlap-free & \tikzcell{	
			\node[normal] (A1) {$u$};
			\node[overlap, right=0pt of A1] (A2) {$v$};
			
			\node[overlap, below=-0.5pt of A2] (B1) {$v$};
			\node[normal, right=0pt of B1] (B2) {$w$};            
		} \\
		\midrule
		
		not infix-free & \tikzcell{
			\node[draw, minimum height=6mm, minimum width=4cm] (u) {$u$};
			\node[normal, below=0pt of u] (v) {$v$};
		} \\
		\midrule
		
		not prefix-free & \tikzcell{        
			\node[draw, minimum height=6mm, minimum width=4cm] (u) {$u$};
			\node[normal, anchor=north west] (v) at (u.south west) {$v$};        
		} \\
		\midrule
		
		not suffix-free & \tikzcell{
			\node[draw, minimum height=6mm, minimum width=4cm] (u) {$u$};
			\node[normal, anchor=north east] (v) at (u.south east) {$v$};        
		} \\
		\midrule
		
		not comma-free & \tikzcell{
			\node[normal] (u) {$u$};
			\node[normal, right=0pt of u] (v) {$v$};
			
			\node[normal, anchor=north] (w) at ($(u.south)!0.5!(v.south)$) {$w$};        
		} \\
		\bottomrule
	\end{tabular}
\end{table}


By definition a code is a subset of $\Sigma^+$, but as $\varepsilon = \varepsilon \cdot \varepsilon$
a language fulfilling \crefitem{constrsync:def:codes}{constrsync:item:code}  
never contains the empty word and we could 
replace $\Sigma^+$ by~$\Sigma^*$  in the definition.
%equivalently define a code as a subset of $\Sigma^*$. 
However, the defining conditions 
for overlap-free languages, infix, prefix, suffix, bifix, solid, and comma-free codes
are also fulfilled for $L = \{\varepsilon\}$ and the exclusion of the empty word makes
a difference in these cases.

\begin{example} \label{constrsync:ex:binarysolidcode}
	\begin{enumerate}
		\item The set $\Sigma$ is a code, as $\Sigma^*$ fulfills the unique factorization property w.r.t.~$\Sigma$.
		
		\item For $n \ge 1$, the set $\Sigma^n$ is a code. A subset of~$\Sigma^n$
		is a solid code if and only if it is overlap-free. Thus, $\Sigma$
		is a solid code.
		
		\item \label{constrsync:ex:comma-free-single-word}  $\{ w \}$ is a code for $w \in \Sigma^+$. 
		Here, we have that 
		$\{ w \}$ is not a comma-free code if and only if $w = uv$
		for $u,v \in \Sigma^*$ such that $uv = vu$.
		This follows since $uv = vu$ implies $uvuv \in \Sigma^+ vu \Sigma^+$
		and, conversely, $ww = x w y$ for $x,y \in \Sigma^+$
		implies that we can write $w = xu = vy$ for $u,v \in \Sigma^+$
		with $uv = w$. Now $xu = uv$ implies $|x| = |v|$
		and similarly we obtain $|u| = |y|$. Thus, since $x$ and $v$
		are equal-length prefixes of~$w$ we have $x = v$
		and similarly $u = y$ and thus $vu = w = uv$.
		
		\item $\{a,ab\}$ is an overlap-free code that is not a prefix code.
		
		\item $\{a,ab,aab\}$ is overlap-free, but not a code since $(aab) = (a)(ab)$.
		
		\item $\{abc,b\}$ is an overlap-free bifix code that is not an infix code.
		
		\item $a^*b$ is a prefix code that is overlap-free and not a suffix code. 
		
		\item  \label{constrsync:ex:binarysolidcodeitem}
		%solid code
		%abab
		%aaabbb
		%aabab
		$\{ aba^ibb \mid i \ge 1\}$ is a solid code. 
		
		\item The language~$C$ from \cref{prelim:lem:alphabet-reduction} forms a solid code.
		
		\item \label{constrsync:ex:comma-free-not-solid} $\{abb,aba\}$ is a comma-free code, but not a solid code since
	    it is not overlap-free.
		
		\item \label{constrsync:ex:infix-not-comma-free} $\{ab,ba\}$ is an infix code, but not a comma-free code.
		
		\item Let $C \subseteq \Sigma^+$ be a prefix code. Then $C\Sigma^n$ is a prefix code for each $n \ge 1$.
			
		\item For $L \subseteq \Sigma^*$ with $L \cap \Sigma^+ \ne \emptyset$,
		the language $L \setminus ( L\Sigma^+ \cup \{\varepsilon\} )$ is a prefix code.
		
		\item \label{constrsync:ex:constr-solid-code} Suppose $C \subseteq \Sigma^+$ is a finite prefix code
		such that $c\Sigma^* \cap C = C \cap \Sigma^*c = \emptyset$
		for some $c \in \Sigma$.
		Set $k = \max\{ l \ge 0 \mid \Sigma^* c^l \Sigma^* \cap C \ne \emptyset \} + 1$.
		For example, when applying the construction to the prefix code $C = \{bcb, bccb\}$
		we obtain $\{cccbcb, cccbccb\}$.
		Then, $c^k C$ is a solid code  as we show next.
		
		\underline{1. $c^kC$ is an infix code:} Assume $xc^kuy = c^kv$ with $u,v \in C$ and $x,y \in \Sigma^*$.
		As $c^k$ is not a factor of $v$, 
		we have $|x| < k$, which gives $x \in c^*$.
		Since $v \notin c\Sigma^*$
		we conclude $|x| = 0$.
		Hence we have $uy = v$ and as $C$ is a prefix code we get $y = \varepsilon$.
		So $c^kC$ is an infix code. 
		
		
		\underline{2. $c^kC$ is overlap-free:} Assume $xy = c^ku$ and $yz = c^kv$
		for $u,v \in C$ and $x,y,z \in \Sigma^*$.
		%with $x,z \ne \varepsilon$.
		For $|y| \le k$, as $yz = c^kv$ we have $y \in c^*$
		and, as $u \notin \Sigma^*c$, from $xy = c^ku$ we get $|y| = 0$.
		For $|y| > k$, from $yz = c^kv$ we find $y = c^kw$ for some $w \in \pref{v}$.
		Then, from $xy = c^ku$, as $c^k \notin \factor{u}$, we must have $|x| < k$.
		Further, as $u \notin c\Sigma^*$, $xy = c^ku$ then gives $|x| = 0$, which 
		implies that $c^ku$ is a prefix of~$c^kv$.
		As $c^kC$ is an infix code, we find $|z| = 0$ and so $y = c^ku = c^kv$.
		Hence $c^kC$ is overlap-free.
		% Assume $y \ne \varepsilon$.
		% As $u \notin \Sigma^*c$, we have $y = c^kv'$ for some $v' \in \Sigma^+$.
		% Then, as $u \notin c\Sigma^*$ we have $|x| > k$.
		% However, this implies $c^k \in \fact{u}$, which 
		% is not possible. Hence, we have $y = \varepsilon$
		% and so~$c^kC$ is overlap-free.
		
		\item By the previous item, for a finite prefix code $C \subseteq \Sigma^+$
		and distinct $c,a \in \Sigma$ with 
		$k = \max\{ l \ge 0 \mid \Sigma^* c^l \Sigma^* \cap C \ne \emptyset \} + 1$
		we have that $c^k a C a$ is a solid code
		and $c^k aC$ and $c^k Ca$ are prefix codes
		if $C \cap \Sigma^*c = \emptyset$ and $c \Sigma^* \cap C = \emptyset$, respectively,
		since $aC$, $Ca$ and $aCa$ are prefix code as well. 
		
	\end{enumerate}
\end{example} 

\begin{remark}		
	We actually came up with the construction from \cref{constrsync:ex:constr-solid-code}
	in \cref{constrsync:ex:binarysolidcode} by analysing
	a result from~\cite[Thm.~29]{FernauGHHVW19}
	and extending it in~\cite[Thm.~17]{DBLP:conf/fct/Hoffmann21} (see also the discussion on page~\pageref{discussion:strongly-self-sync-codes}).	
	However, a closely related variant of this construction is much older and was introduced 
	by Gilbert~\cite{Gilbert1960Synchronization} in the context of synchronization between sender and receiver during message transmission
	and by Levenshtein~\cite{Lev64,levenshtein1970maximum_en} 
	to construct maximal solid codes of equal-length words.\footnote{Levenshtein referred to them as non-overlapping codes. In fact, 
		codes of equal-length words
		are infix codes. Thus, such a code forms a solid code if and only if it is overlap-free.}
	Levenshtein considers codes of the form $c^k C$, where
	$C = \{ u \in \Sigma^{n} \mid u \notin \Sigma^* c^k \Sigma^* \text{ and } u \notin c\Sigma^* \text{ and } u \notin \Sigma^* c \}$
	for $1 \le k < n$ and $c \in \Sigma$.
	Further generalizations appear in~\cite[Construction~2]{Blackburn2015NonOverlappingCodes}. It was also independently rediscovered in~\cite{CheeKiahPurkayasthaWang2013}.
\end{remark}

A prefix code is also characterized by the property
that for code words $u,v$ and any words $w,w'$
the equality $uw = vw$ implies $u = v$
since $uv = vw$ implies that either $u$ is a prefix of~$v$
or vice versa by the unique factorization property of~$\Sigma^*$
w.r.t.~$\Sigma$ and so both must be equal.
Now, as the naming anticipates, this yields that prefix 
codes (and similarly suffix codes) are indeed codes: The equation $u_1 \cdots u_n = v_1 \cdots v_m$ implies
$u_1 = v_1$ and $u_2 \cdots u_n = v_2 \cdots v_m$ for prefix codes (resp.~$u_m = v_m$ and $u_1 \cdots u_{n-1} = v_1 \cdots v_{m-1}$ for suffix codes).
Continuing, we ultimately find $n = m$ and $u_i = v_i$ 
for all $i \in \{1,\ldots,n\}$. 
Thus, as subclasses, bifix and infix codes are also codes.
Furthermore, as we show next, solid and comma-free codes are infix codes,
yielding that all the mentioned classes except overlap-free languages
are codes. See \cref{fig:code-hasse} for an overview of the inclusion relations.

\begin{proposition}[\cite{DBLP:reference/hfl/JurgensenK97,DBLP:journals/tcs/JurgensenS21}]
	\label{constr_sync:prop:selfsync_implies_avoiding}
	\begin{enumerate} 
		\item \label{constr_sync:prop:item:sc-in-cf} The class of solid codes is properly contained in the 
		class of comma-free codes. 
		
		\item \label{constr_sync:prop:item:cf-in-ic} The class of comma-free codes is properly contained in the class of infix codes.
	\end{enumerate} 	
\end{proposition}
\begin{proof}
	A comma-free code that is not a solid code
	and an infix code that is not comma-free
	are given in \cref{constrsync:ex:comma-free-not-solid} and \cref{constrsync:ex:infix-not-comma-free} 
	of \cref{constrsync:ex:binarysolidcode}, respectively.
	Next, we show the claimed inclusions.	
		
	\ref{constr_sync:prop:item:sc-in-cf}:
	Let $L \subseteq \Sigma^+$ be non-empty. Suppose~$L$ is not comma-free.
	Then there exist $u,v,w \in L$ such that
    $uv = xwy$ for some $x,y \in \Sigma^+$.
	Then $|x| \ge |u|$ or $|y| \ge |v|$ 
	implies $w \in \factor{v}$ or $w \in \factor{v}$, respectively, so that $L$
	is not an infix code.
	For $|x| < |u|$ and $|y| < |v|$, we have 
	$w = u'v'$ for a proper non-empty suffix $u'$ of $u$
	and a proper non-empty prefix~$v'$ of~$v$,
	so that $L$ is not overlap-free.
	Hence, if $L$ is a solid code, we must have $L^2 \cap \Sigma^+L\Sigma^+ = \emptyset$.
	
	\ref{constr_sync:prop:item:cf-in-ic}: Let $C \subseteq \Sigma^+$ be a comma-free code.
	Suppose $u, vuw \in C$ for $u,v,w \in \Sigma^*$.
	For $w \ne \varepsilon$, we have $uvuw \in C^2 \cap \Sigma^+ u \Sigma^+$,
	and for $v \ne \varepsilon$, we have $vuwu \in C^2 \cap \Sigma^+u \Sigma^+$.
	Whence, by comma-freeness, we find $v = w = \varepsilon$. So $C$ is an infix code.
\end{proof}

\begin{figure}[ht]
	\centering
	\begin{tikzpicture}[
		node distance=1.7em,
		every node/.style={align=center},
		every path/.style={-}
		]
		
		\node (solid) {solid};
		\node (comma) [above=of solid] {comma-free};
		\node (infix) [above=of comma] {infix};
		\node (bifix) [above=of infix] {bifix};
		\node (prefix) [above left=of bifix] {prefix};
		\node (suffix) [above right=of bifix] {suffix};
		\node (general) [above=of bifix, yshift=10mm] {code};
		
		\node (overlap) [
		right=3em of comma,
		yshift=6mm
		] {overlap-free};
		
		\draw (solid) -- (comma);
		\draw (comma) -- (infix);
		\draw (infix) -- (bifix);
		
		\draw (bifix) -- (prefix);
		\draw (bifix) -- (suffix);
		
		\draw (prefix) -- (general);
		\draw (suffix) -- (general);
		
		\draw (solid) -- (overlap);		
	\end{tikzpicture}
	\caption[Relations between different code classes and overlap-free languages]{Inclusion relations between different code classes and overlap-free languages.}
	\label{fig:code-hasse}
\end{figure}

\IN{rooted tree} \index{root|see{rooted tree}}
\begin{remark} \label{syncmax:rem:root-tree-prefix-code}
	Prefix codes can be visualized as rooted trees.
	A \emph{rooted tree} is an undirected connected and acyclic graph
	with the selection of a distinguished node, called the \emph{root}.
	A prefix code $C \subseteq \Sigma^+$
	induces the rooted tree $(C,  \{ \{ u, ux \} \mid u,ux \in \pref{C} \})$ with root $\varepsilon$ and 
	where each node has outdegree at most~$|\Sigma|$.
	For example, $\{00,01,10,110,111\}$ from \Cref{fig:intro:encoding}
	is a prefix code with its corresponding
	rooted tree shown (with added directions).
	Conversely, by labeling the outgoing edges 
	of each node of a rooted tree where each node has outdegree at most~$|\Sigma|$
	by a distinct letter from~$\Sigma$
	we uniquely obtain a prefix code.
	In this sense, up to a reordering of the alphabet,
	prefix codes over~$\Sigma$
	correspond uniquely to rooted trees with outdegree at most $|\Sigma|$.
\end{remark}

\IN{intercode of index one}
Comma-free codes are also known as \emph{intercodes of index one}~\cite[Thm.~3]{DBLP:journals/iandc/Jurgensen08}.\footnote{There appear to be two typographical errors in~\cite{DBLP:journals/iandc/Jurgensen08}. The definition of intercodes of index one omits the exponent~$2$, i.e., it states $L \cap \Sigma^+ L \Sigma^+ = \emptyset$ instead of $L^2 \cap \Sigma^+ L \Sigma^+ = \emptyset$. Furthermore, the condition $u,w \ne \varepsilon$ appears to be missing from the definition of an overlap-free code. Without this condition, only $L=\emptyset$ satisfies the definition, for if $v \in L$, we have $\varepsilon \cdot v = v \cdot \varepsilon$, and so $v=\varepsilon$. But this is impossible because $L \subseteq \Sigma^+$.}
Originally, comma-free codes were defined to be subsets of $\Sigma^n$ for some $n \ge 1$~\cite{golomb_gordon_welch_1958}, but this restriction was subsequently dropped~\cite{DBLP:books/daglib/0025093,DBLP:journals/iandc/Jurgensen08,DBLP:journals/tcs/JurgensenS21}.

\IN{coding morphism}
\IN{channel alphabet}
\IN{source alphabet}
The (algebraic) approach regarding %(variable-length) 
codes as certain subsets of $\Sigma^*$ goes back to Schützenberger~\cite{DBLP:journals/tit/Schutzenberger56}.
Prior to this, the theory of codes had its origin in Shannon's theory of information~\cite{Shannon:1948:MTCa,Shannon:1948:MTCb}.
Both viewpoints are linked via \emph{coding morphisms}, i.e., injective morphisms $\varphi \colon \Gamma^* \to \Sigma^*$
such that $\applyfunc{\Gamma}{\varphi}$ is a code.
Here, $\Gamma$ is called the \emph{source alphabet} 
and $\Sigma$ the \emph{channel alphabet}. 
Intuitively, this  expresses the fact
that codes  may be used as communication tools via a correspondence (encoding)
between the letters  $a \in \Gamma$ of the source alphabet
by code words $\varphi(a)$ over the channel alphabet~$\Sigma$,
where injectivity guarantees that a message $\varphi(u)$, $u \in \Gamma^*$,
can be uniquely deciphered as $u$ by the receiver.
%the decoding of a will never lead to ambiguities. 
% used for transmission
% of a word over $\Gamma$.
%The unique factorization property 
%guarantees that the decoding will never lead to ambiguities. 
%The link is as follows: codes may be used as communication tools by defining a correspondence (encoding)
%between the letters of a ``source'' alphabet $\Gamma$ 
%and the words of a code $C \subseteq \Sigma^+$ over the ``channel'' alphabet~$\Sigma$. The unique factorization property 
%guarantees that the decoding will never lead to ambiguities. 
%This is formalized by \emph{coding morphisms}: injective morphisms $\varphi \colon \Gamma^* \to \Sigma^*$ such that $\applyfunc{\Gamma}{\varphi}$ is a code where $\Gamma$ is called the \emph{source alphabet} 
%and $\Sigma$ the \emph{channel alphabet}.
For instance, in the introductory
example on error-prone communication from \cref{intro:sec:examples:communication}, 
the source alphabet is $\{a,h,l,o,r\}$
and the channel alphabet is $\{0,1\}$.
%For reference we state the following relation between coding morphisms and codes~\cite[Prop~2.1.1]{DBLP:books/daglib/0025093}.
The following statement~\cite[Prop~2.1.1]{DBLP:books/daglib/0025093} summarizes the relation of codes to coding morphisms.

\begin{proposition}[{\cite[Prop.~2.1.1]{DBLP:books/daglib/0025093}}]
	\label{constr_sync:prop:coding_morphisms}
	To every code $C \subseteq \Sigma^+$ there exists a coding morphism $\varphi \colon \Gamma^* \to \Sigma^*$ with $\varphi(\Gamma) = C$.  
	Moreover,
	any bijection from some alphabet~$\Gamma$ (here possibly infinite) onto a code $C \subseteq \Sigma^+$ 
	extends (as in \cref{prelim:prop:UMP}) to a unique injective morphism from $\Gamma^*$ to $\Sigma^*$ with image~$C^*$. Conversely, for injective morphisms $\varphi \colon \Gamma^* \to \Sigma^*$, $\applyfunc{\Gamma}{\varphi}$
	is a code.	
	%	 = C$ for $C \subseteq \Sigma^+$, then~$C$ is a code.\todo{originalformulierung?}
\end{proposition}

%\todo{großschreiben Universal...?}
%By the universal mapping property (\cref{prelim:prop:UMP}) the extension 
%in \cref{constr_sync:prop:coding_morphisms} 
%is unique and we always have $\varphi(\Gamma^*) = C^*$.

Throughout this work, our alphabet is always finite,
the only exception being \cref{prelim:prop:UMP,constr_sync:prop:coding_morphisms} 
where the (source) alphabets can be infinite.

%Note that the source alphabet might be infinite in this context.

\begin{remark}
	Combining these notions, we find that $L \subseteq \Sigma^+$
	is a code if and only if the map $\hat{i} \colon L^* \to \Sigma^*$
	from \cref{prelim:sec:formallang:basics} is injective, i.e.,
	a coding morphism where~$L$ is considered as a (possibly infinite)
	alphabet.
	% \hat{i} factors through any Gamma* to \Sigma* where Gamma and L bijective. Damit folgt vorherige Prop.
	Thus we can identify words formed by concatenating words from a code~$C$ with tuples
	over~$C$ in a unique way. This yields that \cref{prelim:prop:UMP}
	transfers to codes, i.e., every function from $C$ into a monoid~$M$
	extends uniquely into a monoid morphism from~$C^*$ into~$M$.	
\end{remark}

The next two results are well-known in combinatorics on words. \Cref{constr_sync:lem:two-words-code}
characterizes codes of size two and \cref{constrsync:prop:conjugate-words} gives a factorization result for triples of words $u,v,w$
such that $uv = wu$. 
For a proof of the former result see~\cite[Cor.~1.2.6, Prop.~1.3.1 \& 1.3.2]{lothaire_1997}
or~\cite[Lem.~5.1]{Ginsburg1964BOUNDEDAL} or~\cite[Lem.~3 \& Cor.~4.1]{LyndonSchutzenberger1962} for the latter three equivalences,
and for the latter result see~\cite[Prop.~1.3.4]{lothaire_1997}
or~\cite[Lem.~2]{LyndonSchutzenberger1962}.

\begin{proposition}
	\label{constr_sync:lem:two-words-code}
	\label{constr_sync:lem:common_power}
	Let $u,v \in \Sigma^*$ be distinct. Then the following are equivalent.
	\begin{enumerate}
		\item $\{u,v\}$ is not a code.
		
		\item $u,v \in w^*$ for some $w^*$.
		
		\item $u^i = v^j$ for some $i,j \ge 0$.
		
		\item $uv = vu$.
	\end{enumerate}	
\end{proposition}

\begin{remark}
	By \cref{constr_sync:lem:common_power}, with \cref{constrsync:ex:comma-free-single-word} from \cref{constrsync:ex:binarysolidcode}, 
	$\{w\}$ for $w \in \Sigma^+$ is comma-free if and only if $w = v^i$ for $i \ge 1$
	implies $i = 1$, 
\end{remark}

% Beweis von einem Teil:
% uv = vu
% u = u_1 .... u_n
% Angenommen u_1 = u_{n - m + 1}, ..., u_m = u_n und u_{m+1} = u_1, ..., u_n = u_{n - m} für ein m.
% Dann potenz von u_1 ... u_m. Folgt ganz einfach, da man die "Gleichungen" einfach nacheinander liest
\begin{proposition} \label{constrsync:prop:conjugate-words}
	\label{constrsync:lem:schuetz-lyndon}
	Let $u,v,w \in \Sigma^*$
	with $w \ne \varepsilon$ such that $uv = wu$.
	Then $w = xy$, $u = (xy)^kx$, and $v = yx$
	for some $x,y \in \Sigma^*$ and $k \ge 0$.		
\end{proposition}

\subsection{More on Prefix Codes}
\label{constrsync:sec:more-on-prefix-codes}

Here we give further properties of prefix codes. We begin with another characterization in \cref{constrsync:prop:char-prefix-code}. \Cref{constrsync:prop:suffix-factorizations} gives additional conditions equivalent to part of this characterization. As pointed out in \cref{constrsync:rem:prefix-codes-general-factorization}, these statements characterize prefix codes in a way that generalizes the defining property of codes by including certain suffixes in addition to code words. We then derive several consequences that will be needed later. Finally, in \cref{constr_sync:prop:aut_prefix_free}, we characterize regular prefix codes via PDFAs.

\begin{proposition} \label{constrsync:prop:char-prefix-code}
	A language $C \subseteq \Sigma^+$ is a prefix code if and only if it is a code
	such that $C^* \cap C^*u \ne \emptyset$ for $u \in \Sigma^*$
	implies $u \in C^*$.
\end{proposition}
\begin{proof}
 First, suppose~$C$ is a prefix code.
 Then it is a code.
 The remaining claim is equivalent to the fact that $vu \in C^*$ for $v \in C^*$
 implies $u \in C^*$.
 We show this by induction on the length of a product of~$v$ into words from~$C$.
 For $v = \varepsilon$, the claim clearly holds.
 Otherwise, we have $v = xx'$ and $vu = yy'$ with $x,y \in C$
 and $x',y' \in C^*$.
 Since~$C$ is a prefix code and $yy' = vu = xx'u$,
 we find $y = x$ and $y' = x'u$.
 Since $y',x' \in C^*$, inductively, we find $u \in C^*$.
 Thus, by induction, we obtain the claim. 
 
 Next, suppose that $C$ is not a prefix code. Then there exist $u \in C$ and a non-empty word $v \in \Sigma^+$ such that $uv \in C$. Thus either $C$ is not a code, or, if $C$ is a code, necessarily $v \notin C^+$. Indeed, otherwise the word $uv \in C$ would have the alternative factorization
 $
 uv = u \cdot v
 $
with $u \in C$ and $v \in C^+$, contradicting the unique factorization property of a code.
\end{proof}

The latter condition in \cref{constrsync:prop:char-prefix-code} is equivalent to two other conditions concerning the uniqueness of certain suffixes, as stated below. Notably, these equivalences hold for arbitrary languages in $\Sigma^+$, 
not just codes. For example, they apply to~$(\Sigma^i)^+$ for $i \ge 1$ and $\{ab,b\}$, although neither of these languages is a code.

\begin{proposition}  \label{constrsync:prop:suffix-factorizations}
	Let $L \subseteq \Sigma^+$ be any language not containing the empty word. Then the following are equivalent.
	\begin{enumerate}
		\item \label{constrsync:prop:suffix-factorizations:1} For $u,v \in \Sigma^* \setminus L\Sigma^*$ with $L^* u \cap L^* v \ne \emptyset$ we have $u = v$. 
		
		\item \label{constrsync:prop:suffix-factorizations:2} $L^* \cap L^* v \ne \emptyset$ for $v \in \Sigma^*$ implies $v \in L^*$.
		
		\item \label{constrsync:prop:suffix-factorizations:3} For $u,v \in \Sigma^+ \setminus L\Sigma^+$ with $L^* u \cap L^* v \ne \emptyset$ we have $u = v$.
	\end{enumerate}
\end{proposition}
\begin{proof}
	$\ref{constrsync:prop:suffix-factorizations:1} \Rightarrow \ref{constrsync:prop:suffix-factorizations:2}$: 
	Let $u, uv \in L^*$. Since $\varepsilon \notin L$, we can write $v = xy$
	with $x \in L^*$ and $y \notin L\Sigma^*$.
	Then we have $uv \in L^* \cap L^*y$. By \cref{constrsync:prop:suffix-factorizations:1}, we find $y = \varepsilon$
	and thus $v \in L^*$.
	
	$\ref{constrsync:prop:suffix-factorizations:2} \Rightarrow \ref{constrsync:prop:suffix-factorizations:3}$: 
	Let $w,w' \in L^*$ such that $wu = w'v$ with $u,v \in \Sigma^+ \setminus L\Sigma^+$.
	Assume $u \ne v$. Then, without loss of generality,
	we have $u = xv$ for some non-empty $x \in \Sigma^+$.
	This gives $w' = wx$. By assumption, we find $x \in L^+$
	and thus $u \in L\Sigma^+$, but $u \notin L \Sigma^+$ by choice of~$u$.
	Thus, we must have $u = v$.
	
	$\ref{constrsync:prop:suffix-factorizations:3} \Rightarrow \ref{constrsync:prop:suffix-factorizations:1}$: Suppose \cref{constrsync:prop:suffix-factorizations:3} holds true.
	Let $L^* u \cap L^* v \ne \emptyset$ for some $u,v \notin L\Sigma^*$.
	Then there exists $w \in L^*$
	such that $wu \in L^*v$.
	Assume $u \ne v$. 
	Then, by \cref{constrsync:prop:suffix-factorizations:3},
	we have $u,v \notin \Sigma^+ \setminus L \Sigma^+$.
	Since $u,v \notin L \Sigma^+$,
	this yields $u = \varepsilon$ or $v = \varepsilon$.
	Without loss of generality, let $u = \varepsilon$.	
	Write $w = xx'$ and $x'x' = yy'$ with $x$ and $y$ being the longest proper prefixes of~$w$ and $x'x'$, respectively, in~$L^*$. Then we have 
	\[
	x', y' \in \Sigma^+ \setminus L\Sigma^+
	\text{ and }
	wx' \in L^* x' \cap L^* y'.
	\]
	By assumption, this gives $x' = y'$.
	Since $x'x' = yy'$ we find $x' = y$.
	Thus $x' \in L^*$ since $y \in L^*$. With $x' \notin L\Sigma^+$ and $|x| > 0$,
	we get $x' \in L$.	
	We have $w \in L^* x' \cap L^* v$ where $v \in \Sigma^+$ since $u \ne v$.
	By assumption, we find $x' = v$, contradicting $v \notin L \Sigma^*$ since $x' \in L$.
	Thus, we must have $v = \varepsilon$.
\end{proof}

Note that $u$ and~$v$ must be non-empty in \cref{constrsync:prop:suffix-factorizations:3} of
\cref{constrsync:prop:suffix-factorizations}.
For example, consider any non-empty language $L \subseteq \Sigma^+$.
Then for $w \in L$ we have $w \in L^*\varepsilon \cap L^*w$,
but $w \ne \varepsilon$ since~$L$ is a subset of~$\Sigma^+$.

Combining \cref{constrsync:prop:char-prefix-code,constrsync:prop:suffix-factorizations},
we obtain that prefix codes extend the unique factorization property of codes
in the following sense.

\begin{corollary}  \label{constrsync:rem:prefix-codes-general-factorization}
	 \label{contrsync:prop:prefix-code-suffix-unique}
	A language $C \subseteq \Sigma^+$ is a prefix code if and only if
	for $u_1, \ldots, u_{n-1} \in C$, $v_1, \ldots, v_{m-1} \in C$,
	and $u_n, v_m \in \Sigma^+ \setminus C\Sigma^+$
	such that
	\[
	u_1 \cdots u_{n-1} u_n = v_1 \cdots v_{m-1} v_m
	\]
	we have $n = m$ and $u_i = v_i$ for all $i \in \{1,\ldots,n\}$.
\end{corollary}
\begin{proof}
	Suppose~$C$ is a prefix code and let $u_1, \ldots, u_n$
	and $v_1, \ldots, v_m$ as in the statement
	with $u_1 \cdots u_{n-1} u_n = v_1 \cdots v_{m-1} v_m$.
	Thus $C^* u_n \cap C^* v_m \ne \emptyset$.
	By \cref{constrsync:prop:char-prefix-code,constrsync:prop:suffix-factorizations}
	this gives $u_n = v_m$.
	Thus $u_1 \cdots u_{n-1} = v_1 \cdots v_{m-1}$.
	Since $C$ is a code, we find $n = m$ and $u_1 = v_1, \ldots, u_{n-1} = v_{n-1}$.	
	
	Conversely, this condition yields that~$C$ is a code
	since it generalizes the unique factorization property.
	Next, suppose $C^* u \cap C^* v$ for some $u,v \in \Sigma^+ \setminus C\Sigma^+$.
	Then there exist $u_1, \ldots, u_n \in C$
	and $v_1, \ldots, v_m \in C$
	such that $u_1 \cdots u_n u = v_1 \cdots v_m v$.
	By assumption, we find $n = m$, $u_i = v_i$ for $i \in \{1,\ldots,n\}$,
	and $u = v$.
	Thus, by \cref{constrsync:prop:char-prefix-code,constrsync:prop:suffix-factorizations},
	we conclude that~$C$ is a prefix code.
\end{proof}

Note that for a prefix code, we have
$C \subseteq \Sigma^+ \setminus C\Sigma^+$, i.e., the condition from
\cref{constrsync:rem:prefix-codes-general-factorization}
also allows $u_n$ and $v_m$ to belong to~$C$.

With our result, a variant of
\cref{constrsync:rem:prefix-codes-general-factorization}
can also be formulated with
$u_n, v_m \in \Sigma^* \setminus C\Sigma^*$
instead of
$u_n, v_m \in \Sigma^+ \setminus C\Sigma^+$.

\begin{corollary}
	Let $C \subseteq \Sigma^+$ be a prefix code.
	Then the factorization witnessing a word to
	be in $\pref{C^*} = C^*( \pref{C} \setminus C )$
	is unique, i.e., 
	for $u \in \pref{C^*}$
	there exists precisely one factorization
	of the form $u = u_1 \cdots u_{n-1} u_n$
	with $u_1, \ldots, u_{n-1} \in C$
	and $u_n \in \pref{C}\setminus C$.	
\end{corollary}
\begin{proof}
	This follows directly from \cref{contrsync:prop:prefix-code-suffix-unique}
	since $\pref{C} \cap C\Sigma^+ = \emptyset$ for a prefix code~$C$.
\end{proof}





\begin{corollary} \label{contrsync:lem:prefix-code-uv}
	Let $C \subseteq \Sigma^+$ be a prefix code.
	Then $\pref{C^*} \cap C^*u \ne \emptyset$ for $u \in \Sigma^*$
	implies $u \in \pref{C^*}$.	
\end{corollary}
\begin{proof}
	Let $vu \in \pref{C^*}$ with $v \in C^*$.	
	Write $vu = xx'$ and $u = yy'$ with $x \in C^*$ and $x'$ being a proper prefix of some word in~$C$
	and $y$ being a longest proper prefix of~$u$ in~$C^*$.
	Since~$C$ is a prefix code, we have $x' \notin C\Sigma^+$
	and, by choice of~$y$, we have $y' \in \Sigma^+ \setminus C\Sigma^+$.
	For $x' = \varepsilon$, we have $vu \in C^*$
	and by \cref{constrsync:prop:char-prefix-code} we find $u \in C^*$.
	Next, suppose $x' \ne \varepsilon$.
	We have $xx' = vu = vyy'$ and so $vu \in C^*x' \cap C^*y'$
	and by \cref{constrsync:prop:char-prefix-code,constrsync:prop:suffix-factorizations} we find $x' = y'$.
	Since $x' \in \pref{C^*}$,
	this gives $u = yx' \in C^*\pref{C^*} = \pref{C^*}$.	
\end{proof}

\Cref{contrsync:lem:prefix-code-uv} fails for general codes. For example, for $C = \{a,ab\}$
we have $ab \in \pref{C^*}$, $a \in C$, but $b \notin \pref{C^*}$.

\begin{corollary}
	Let $C \subseteq \Sigma^+$ be a prefix code and $u \in \Sigma^*$.
	Suppose $C^*u \cap C^* ( \pref{C} \setminus C ) \ne \emptyset$
	and that~$u$ is not a proper prefix of some word from~$C$ and $u \notin C\Sigma^+$.
	Then we have $u \in C$.
\end{corollary}
\begin{proof}
	Since $C^*w \subseteq \pref{C^*}$,
	we have $C^* u \cap \pref{C^*} \ne \emptyset$.
	Thus, by \cref{contrsync:lem:prefix-code-uv}, 
	we find $u\in \pref{C^*}$.
	Since $u \notin C\Sigma^+$, this gives $u \in \pref{C}$.
	Then, since $u \notin \pref{C}\setminus C$,
	we must have $u \in C$.
\end{proof}

\begin{corollary}
	Let $C \subseteq \Sigma^+$ be a prefix code and $L \subseteq \Sigma^*$
	such that $C^* L \subseteq L$.
	Then we have the following:
	\begin{align}  
		C^* ( L \setminus \pref{C^*} ) & = L \setminus \pref{C^*} \label{constrsync:eqn:complementprefCstar} \\	
		( \pref{C} \setminus C ) \Sigma	\cap \pref{C^*} & = ( \pref{C} \setminus C ) \Sigma	\cap \pref{C}, \label{constrsync:eqn:capprefCstar} 
	\end{align}
\end{corollary}
\begin{proof}
	\Cref{constrsync:eqn:complementprefCstar}: Since $\varepsilon \in C^*$
	we have $L \setminus \pref{C^*} \subseteq 	C^* ( L \setminus \pref{C^*} )$.
	Next, consider $u \in C^*$
	and $v \in L \setminus \pref{C^*}$.
	Now, \cref{contrsync:lem:prefix-code-uv} 
	would yield $v \in \pref{C^*}$
	for $uv \in \pref{C^*}$.
	Thus, we must have $uv \in  L \setminus \pref{C^*}$.
	
	\Cref{constrsync:eqn:capprefCstar}:	
	Let $ux = v_1 \cdots v_n$ for $u \in \pref{C} \setminus C$, $x \in \Sigma$,
	and $v_1, \ldots, v_{n-1} \in C$
	and $v_n \in \pref{C}$.
	We have $ux \notin C\Sigma^+$.
	Thus $n = 1$ and so $ux \in \pref{C}$.
\end{proof}

Lastly, we collect some facts on PDFAs accepting prefix codes. Related statements to Item~\ref{constr_sync:prop:aut_prefix_free:item:trim} and~\ref{constr_sync:prop:aut_prefix_free:item:min}
from \ref{constr_sync:prop:aut_prefix_free} have been given in~\cite[Prop.~2.2]{Eilenberg1974}.

\begin{proposition}
	\label{constr_sync:prop:aut_prefix_free}
	\label{constr_sync:prop:aut_prefix_free_exact_sync}
	Let $L \subseteq \Sigma^+$ be a regular language.
	
	\begin{enumerate}
		\item \label{constr_sync:prop:aut_prefix_free:item:trim} $L$ is a prefix code if and only if for every trim PDFA $\automat{B}{P}{\Sigma}{\mu}{p_0}{E}$
		with $L = L(\mathcal B)$ the states in $E$ have no outgoing transitions, i.e., $\mu(p, x)$ is undefined
		for each $p \in E$ and $x \in \Sigma$.
		
		% det und PDFA equal wenn maximal, sonst determinstic eien non-final trap zustand
		\item  \label{constr_sync:prop:aut_prefix_free:item:min} If $L$ is a prefix code, then the complete and deterministic minimal automaton $\minaut{L}$
		and a PDFA with the least number of states accepting $L$
		both have a single final state.
		
		\item  \label{constr_sync:prop:aut_prefix_free:item:sync} If $L$ is a prefix code and $\automat{B}{P}{\Sigma}{\mu}{p_0}{\{e\}}$ a  trim PDFA 
		with $L = L(\mathcal B)$, 
		then there exists an exactly synchronizing word $w \in L(\mathcal B)$
		with $\mu(P, w) = \{e\}$.
	\end{enumerate}
\end{proposition} 
\begin{proof}
	\ref{constr_sync:prop:aut_prefix_free:item:trim}: Let $\automat{B}{P}{\Sigma}{\mu}{p_0}{E}$ be a trim PDFA with $L = L(\mathcal B)$.
	First, suppose $\mu(p, x)$ is defined for some $p \in E$ and $x \in \Sigma$.
	As $\mathcal B$ is trim, 
	there exist $u,v \in \Sigma^*$
	with $\mu(\mu(p, x), u) \in E$ and $\mu(p_0, v) = p$.
	As $v \in L$ is a proper prefix of $vxu \in L$, $L$ is not a prefix code.
	Conversely, suppose $L$ is not a prefix code and $L = L(\mathcal B)$.
	Then there exist $u,v \in L$
	such that $v = uxu'$ for some $x \in \Sigma$ and $u' \in \Sigma^*$.
	Then $\mu(p_0, u) \in E$ and $\mu(p_0, ux)$ is defined.
	
	\ref{constr_sync:prop:aut_prefix_free:item:min}: 
	For $u,v \in L$ and $w \in \Sigma^+$
	we have $uw, vw \notin L$ since~$L$ is a prefix code,
	which gives $u \equiv_L v$. Therefore we obtain the claim
	for the complete and deterministic minimal automaton.
	The PDFA with the least number of states (up to renaming of the states)
	is obtained from the minimal automaton by omitting
	all non-coaccessible states (which form a set of size at most one), i.e.,
	by omitting the state $[u]_{\equiv_L}$ for $u \notin \pref{L}$
	if such a word $u$ exists. Formally, 
	let $\automat{B}{P}{\Sigma}{\mu}{p_0}{F}$
	where $P = \{ [u]_{\equiv_L} \mid u \in \pref{L} \}$, $p_0 = [\varepsilon]_{\equiv_L}$,
	$F = \{ [u]_{\equiv_L} \mid u \in L \}$,
	and $\mu([u]_{\equiv_L}, x) = [ux]_{\equiv_L}$ for $x \in \Sigma$ and $u \in \pref{L}$. Since~$L$ is regular, $P$ is finite.
	We have $L = L(\mathcal B)$ and $F$ is a singleton, which gives the claim.
		
	\ref{constr_sync:prop:aut_prefix_free:item:sync}: Set $n = |P|$. By Item~\ref{constr_sync:prop:aut_prefix_free:item:trim}
	the state $e$ has no outgoing transitions.
	Inductively, define a sequence of subsets of states $P_0, P_1, P_2, \ldots$, all containing the state $e$, % |P_i| \le n - i für i \ge n - 1
	by setting $P_0 = P$
	and, for $i > 0$, $P_i = \mu(P_{i-1}, u_i)$
	where $u_i$ is some word with $\mu(p, u_i) = e$ if there exists $p \in P_{i-1} \setminus \{e\}$
	and $u_i = \varepsilon$ otherwise.
	Then, for $|P_i| > 1$, as $u_i \ne \varepsilon$ yields that $\mu(e, u_i)$ is undefined,
	it is $|P_{i+1}| < |P_i|$.
	So $\mu(P, u_0 u_1 \cdots u_{n-1}) = \{e\}$ and there exist $p \in P$ and $v \in \Sigma^*$
	such that $\mu(p, u_0 u_1 \cdots u_{n-1}) = e$ and $\mu(p_0, v) = p$.
	Then for $w = vu_0 u_1 \cdots u_{n-1}$ we have $w \in L(\mathcal B)$ and $\mu(P, w) = \{e\}$.				
\end{proof}

\begin{example}\label{constrsync:ex:exact-sync}
	Consider the prefix code $C = \{b,ab,aab\}$.
	Then we obtain a trim PDFA $\mathcal B$ 
	accepting $C$ by setting $\mathcal B = (P, \{a,b\}, \mu, \varepsilon, \{e\})$ with $P = \pref{C}\setminus C \cup \{e\}$, i.e., the proper prefixes
	of~$C$ and~$e$ are considered as states of $\mathcal B$,
	and 
	\[
	\mu(u, x) = \begin{cases}
		e   & \text{if } ux \in C, \\
		ux  & \text{if } ux \in \pref{C} \setminus C.
	\end{cases}
	\]
	We have $\mu(P, b) = \{e\}$ (we also obtain this word by following the procedure from the proof of 
	\crefitem{constr_sync:prop:aut_prefix_free}{constr_sync:prop:aut_prefix_free:item:sync}).
	Other exactly synchronizing words are $aa$ and $aab$ with $\mu(P, aa) = \{aa\}$ and $\mu(P, aab) = \{e\}$.
\end{example}

\subsection{Relations to the Levenshtein-Romanov Mapping}
\label{constrsync:sec:codes:levensthein-romanov}

Here, we state characterizations and conditions for prefix, suffix, bifix, and infix codes
and overlap-free languages and combinations of these properties. We collect
several conditions here, but our focus is on those involving the Levenshtein-Romanov mapping.

\begin{proposition} \label{constrsync:prop:prefix-code}
	Let $L \subseteq \Sigma^+$ be non-empty. Then the following are equivalent.
	\begin{enumerate}
		\item \label{constrsync:prop:prefix-code:1} $L$ is a prefix code.
		
		\item \label{constrsync:prop:prefix-code:1a} $\pref{L} \setminus L$ equals the
		set of proper prefixes $\bigcup_{u \in L} ( \pref{u} \setminus \{u\} )$ of words from~$L$.
		
		\item \label{constrsync:prop:prefix-code:3} We have $\sigma_{\pref{L} \setminus \{u\}}(u) \ne u$ 
		for all $u \in L$.\footnote{With \cref{constrsync:eqn:sigma_prefL}, this is equivalent to $L$ being a singleton or, otherwise, to $\sigma_{L\setminus \{u\}}(u) \ne u$ for all $u \in L$.}
		
		\item \label{constrsync:prop:prefix-code:4} For all $u \in L$ and $v \in \Sigma^+$
		we have $\sigma_L(uv) \ne uv$.	
	\end{enumerate}
\end{proposition}
\begin{proof}
	$\ref{constrsync:prop:prefix-code:1} \Rightarrow \ref{constrsync:prop:prefix-code:1a}$:
	The set $\pref{L} \setminus L$ is always a subset of the set of proper prefixes of words from~$L$.
	Conversely, if~$u \in \pref{L}$ such that $uv \in L$ for some non-empty $v\in \Sigma^+$,
	then $u \notin L$ by the defining condition of a prefix code.	
		
	$\ref{constrsync:prop:prefix-code:1a} \Rightarrow \ref{constrsync:prop:prefix-code:3}$: 
	Let $u \in L$ with $\sigma_{\pref{L}\setminus{u}}(u) = u$.
	Then $u$ is a proper prefix of a word $v \in L$ with $u \ne v$.
	We have $u \notin \pref{L} \setminus L$
	and thus $\pref{L} \setminus L$ does not equal the set of proper prefixes.
		
	$\ref{constrsync:prop:prefix-code:3} \Rightarrow \ref{constrsync:prop:prefix-code:4}$: Suppose $\sigma_L(uv) = uv$ for $u \in L$
	and $v\in \Sigma^+$. Then we have $uv \in \pref{L}$
	where $uv \ne u$. Thus $L$ is not a singleton set 
	and we find $\sigma_{L\setminus\{u\}}(uv) = uv$.
	
	$\ref{constrsync:prop:prefix-code:4} \Rightarrow \ref{constrsync:prop:prefix-code:1}$: 
	If $\sigma_L(uv) = uv$
	for $u \in L$ and $v \in \Sigma^+$, then $uv\in \pref{L}$
	and thus $uvw \in L$ for some $w \in \Sigma^*$.
	Since $vw \ne \varepsilon$, $L$ is not a prefix code in this case.	
\end{proof}

\begin{proposition} \label{constrsync:prop:suffix-code}
	Let $L \subseteq \Sigma^+$ be non-empty. Then the following are equivalent.
	\begin{enumerate}
		\item \label{constrsync:prop:suffix-code:1} $L$ is a suffix code.\footnote{By \crefitem{constrsync:prop:suffsigmaB}{constrsync:prop:suffsigma:3B}, $L$ being a suffix code
		is also equivalent to $\suff{\sigma_{\pref{L}\setminus\{u\}}} \cap ( L \setminus \{u\} ) = \emptyset$
		for all $u \in L$.}
		
		\item \label{constrsync:prop:suffix-code:1a} $\suff{L} \setminus L$ equals the set of proper suffixes
		$\bigcup_{u \in L} ( \suff{u}\setminus\{u\} )$ of words from~$L$.
				
		\item \label{constrsync:prop:suffix-code:4} 
		%We have $uv \notin L$ for all $v \in L$ and $u \in \Sigma^+$.
		We have $\sigma_L(uv) \notin L \setminus\{v\}$ for all $v \in L$ and $u \in \Sigma^+$.
	\end{enumerate}
\end{proposition}
\begin{proof}
	$\ref{constrsync:prop:suffix-code:1} \Rightarrow \ref{constrsync:prop:suffix-code:1a}$: The set $\suff{L}\setminus L$
	is contained in the set of proper suffixes for any language. Conversely, if~$u$ is a proper suffix of a word from~$L$,
	then it is a suffix not contained in~$L$ if~$L$ is a suffix code, i.e., it is contained in $\suff{L}\setminus L$.
			
	$\ref{constrsync:prop:suffix-code:1a} \Rightarrow \ref{constrsync:prop:suffix-code:4}$: 
	Suppose $\sigma_L(uv) \in L \setminus \{v\}$ for some $v \in L$
	and $u \in \Sigma^+$. Then $\sigma_L(uv)$ is a word in~$L$
	having~$v$ as a proper suffix, but we have $v \notin \suff{L} \setminus L$ since $v \in L$.
		
	$\ref{constrsync:prop:suffix-code:4} \Rightarrow \ref{constrsync:prop:suffix-code:1}$: For $uv \in L$ with $v \in L$ and $u \in \Sigma^+$ we have $\sigma_L(uv) = uv \in L \setminus \{ v \}$.
\end{proof}

\begin{example}
	 $L = \{aaab,ab,aabb\}$ is not a suffix code and $\sigma_{L\setminus\{u\}}(u) \notin L$ for all $u \in L$.
%	\begin{enumerate}
%		\item $L = \{baa,aaab,a\}$ is not a suffix code
%		and we have $\sigma_{L\setminus\{u\}}(u) \notin L \setminus \{u\}$
%		for all $u \in L$.
%	\end{enumerate}
\end{example}

\begin{proposition} \label{constrsync:prop:bifix-code}
	Let $L \subseteq \Sigma^+$ be non-empty. Then the following are equivalent.
	\begin{enumerate}
		\item \label{constrsync:prop:bifix-code:1} $L$ is a bifix code.
		
		\item \label{constrsync:prop:bifix-code:2} We have $u,v \notin \suff{u}\cap \pref{v}$ for all $u,v \in L$ with $u\ne v$.
		
		\item \label{constrsync:prop:bifix-code:3} We have $\suff{\sigma_{\pref{L}\setminus\{u\}}(u)} \cap L = \emptyset$	
		for all $u \in L$.
	\end{enumerate}
\end{proposition}
\begin{proof}
$\ref{constrsync:prop:bifix-code:1} \Leftrightarrow \ref{constrsync:prop:bifix-code:2}$: 
For distinct $u,v \in L$ we have $u \in \suff{u} \cap \pref{v}$ (resp.~$v \in \suff{u} \cap \pref{v}$)
if and only if $u$ is a proper prefix of~$v$ (resp.~$v$ is a proper suffix of~$v$) in~$L$,
which gives the claim.

$\ref{constrsync:prop:bifix-code:2} \Leftrightarrow \ref{constrsync:prop:bifix-code:3}$: 
By \crefitem{constrsync:prop:suffsigmaB}{constrsync:prop:suffsigma:3B}, we have $v \in \suff{\sigma_{\pref{L}\setminus\{u\}}(u)} \cap L$ if and only if $v$ is a proper suffix of $u$
or $v = u$ and $u$ is a proper prefix of some word from~$L$,
which gives the claim.
\end{proof}

Clearly, $L \subseteq \Sigma^+$ is a prefix (resp.~suffix) code
if and only if $u \notin \suff{u} \cap \pref{v}$ (resp.~$u \notin \suff{v} \cap \pref{u}$) for distinct $u,v \in L$, 
giving similar characterizations as \cref{{constrsync:prop:bifix-code:2}} of \cref{constrsync:prop:bifix-code}.

\begin{proposition} \label{constrsync:prop:infix-code}
	Let $L \subseteq \Sigma^+$ be non-empty. Then the following are equivalent.
	\begin{enumerate}	
		\item \label{constrsync:prop:infix-code:1} $L$ is an infix code.	
		
		\item \label{constrsync:prop:infix-code:3} $L \Sigma^+ \cap \suff{L} = \Sigma^+ L \cap L = \emptyset$.
		
		\item \label{constrsync:prop:infix-code:4} $\Sigma^* L \Sigma^+ \cap L = \Sigma^+ L \Sigma^* \cap L = \emptyset$.
		
		\item \label{constrsync:prop:infix-code:5} $\Sigma^* L \Sigma^+ \cap L = \Sigma^+ L \cap L = \emptyset$.		
		
		% entsprechende chars wie die folgenden gehen auch für prefix und suffix codes, im fliestext danach
		\item \label{constrsync:prop:infix-code:2} $\Sigma^+ L \cap \pref{L} = L \Sigma^+ \cap L = \emptyset$.
		
		\item \label{constrsync:prop:infix-code:6} $L$ is a prefix code such that $\sigma_L(uv) = v$ for all $u \in \Sigma^*$ and $v \in L$.
			
		\item \label{constrsync:prop:infix-code:8} We have $uv \notin \pref{L}$
		and $vu \notin L$ for all $u \in \Sigma^+$ and $v \in L$.
				
		\item \label{constrsync:prop:infix-code:9} We have $\suff{\sigma_{\pref{L}\setminus\{u\}}(u)} \cap L = \emptyset$ for all $u \in \pref{L}$.
	\end{enumerate}
\end{proposition}
\begin{proof}
The equivalences from \cref{constrsync:prop:infix-code:1,constrsync:prop:infix-code:3,constrsync:prop:infix-code:4,constrsync:prop:infix-code:5,constrsync:prop:infix-code:2} are clear.
	
$\ref{constrsync:prop:infix-code:2}\Rightarrow \ref{constrsync:prop:infix-code:6}$:
The condition $\Sigma^+L \cap L = \emptyset$ implies that~$L$ is a prefix code.
Since $v \in L$ we have $\sigma_L(uv) \in \Sigma^*v \cap \pref{L}$.
By assumption, $\Sigma^+ v \cap \pref{L}$ is empty. Thus $\sigma_L(uv) = v$.

$\ref{constrsync:prop:infix-code:6}\Rightarrow \ref{constrsync:prop:infix-code:8}$:
Let $v \in L$ and $u \in \Sigma^+$.
For $uv \in \pref{L}$, by \cref{constr_sync:eqn_sigma_pref1} we have $\sigma_L(uv) = uv \ne v$.
For $vu \in L$, $L$ is not a prefix code.

$\ref{constrsync:prop:infix-code:8}\Rightarrow \ref{constrsync:prop:infix-code:9}$:
Let $v \in \suff{\sigma_{\pref{L}\setminus\{u\}}(u)} \cap L$.
By \cref{constrsync:prop:suffsigma:3B} of \cref{constrsync:prop:suffsigmaB}, 
$v$ is a proper suffix of $u$ or we have $u = v$ and $u$ is a proper prefix of some word from~$L$.
In the first case, we have $wv \in L$ for some $w \in \Sigma^+$
and in the latter case $vw \in L$ for some $w \in \Sigma^+$.
Thus, \cref{constrsync:prop:infix-code:8} does not hold.

$\ref{constrsync:prop:infix-code:9}\Rightarrow \ref{constrsync:prop:infix-code:1}$:
Let $v \in L$ and $uvw \in L$ with $u \ne \varepsilon$ or $w \ne \varepsilon$.
For $u = \varepsilon$, $v$ is a proper prefix of some word from~$L$.
Thus we have $v = \sigma_{\pref{L}\setminus\{v\}}(v)$ and $v \in \pref{L}$.
So,  \cref{constrsync:prop:infix-code:9}
does not hold.
Next, suppose $u \ne \varepsilon$.
Since $v \in L$ and $uv\ne v$
we have $\sigma_{\pref{L}\setminus\{uv\}}(uv) = w'v$ for some $w' \in \Sigma^*$.
Thus $v \in \suff{\sigma_{\pref{L}\setminus\{uv\}}(uv)} \cap L$.
Since $uv \in \pref{L}$, \cref{constrsync:prop:infix-code:9}
does not hold.
\end{proof}

\begin{example} 	
	For $L = \{a,ab\}$, we have $\sigma_L(uv) = v$ for all $u \in \Sigma^*$ and $v \in L$, but $L$ is not a prefix code. Hence, the assumption that $L$ is a prefix code in \crefitem{constrsync:prop:infix-code}{constrsync:prop:infix-code:6} is necessary.
\end{example}

Similar characterizations as in \cref{constrsync:prop:infix-code:2,constrsync:prop:infix-code:3,constrsync:prop:infix-code:4,constrsync:prop:infix-code:5} 
from \cref{constrsync:prop:infix-code} can also be easily given for prefix and suffix codes.

For $u,v \in \Sigma^*$, $\sigma_{\{v\}}(u)$ is the longest word $w$
such that $u \in w\Sigma^*$ and $v \in \Sigma^* w$. This suggests that we can relate
overlap-freeness to the Levenshtein-Romanov mapping. The next statements give some such relations.

\begin{proposition} \label{constrsync:prop:overlap-free}
	Let $L \subseteq \Sigma^+$ be non-empty. Then the following are equivalent.
	\begin{enumerate}
		\item \label{constrsync:prop:overlap-free:1} $L$ is overlap-free.
		
		\item \label{constrsync:prop:overlap-free:3} For all $u \in \Sigma^+$ we 
		have $u\Sigma^+\cap L = \emptyset$ or $\Sigma^+ u \cap L = \emptyset$.
		
		\item \label{constrsync:prop:overlap-free:2} For all $u,v \in L$ we 
		have $\suff{u} \cap \pref{v} \subseteq \{\varepsilon, u, v\}$.
		
		\item \label{constrsync:prop:overlap-free:2a} For all $u,v \in L$
		we have $\sigma_{\{v\}}(u) \in \{\varepsilon, u, v\}$ and $\suff{u}\cap \pref{u} \subseteq \{\varepsilon, u\}$.	
				
		\item \label{constrsync:prop:overlap-free:2aa} For all $u,v \in L$ we have $\sigma_{\pref{v}\setminus\{v\}}(u) \in \{\varepsilon, u\}$.
		
		\item \label{constrsync:prop:overlap-free:3a} We have $uv \notin L$ for all non-empty and proper prefixes $u,v$ of words from~$L$.
	\end{enumerate}
\end{proposition}
\begin{proof}
$\ref{constrsync:prop:overlap-free:1} \Rightarrow \ref{constrsync:prop:overlap-free:3}$:  Suppose we have $u\Sigma^+ \cap L \ne \emptyset$
and $\Sigma^+u \cap L \ne \emptyset$ for some $u \in \Sigma^+$.
Then $uv \in L$ and $wu \in L$ for some $v,w \in \Sigma^+$ and so~$L$ is not overlap-free.

$\ref{constrsync:prop:overlap-free:3} \Rightarrow \ref{constrsync:prop:overlap-free:2}$: 
Let $w \in \suff{u} \cap \pref{v}$
with $w \notin\{\varepsilon,u,w\}$ for $u,v \in L$.
Then $u = xw$ and $v = wy$ for some $x,y \in \Sigma^+$ and so $u \in \Sigma^+w$ and $v \in w\Sigma^+$.

$\ref{constrsync:prop:overlap-free:2} \Rightarrow \ref{constrsync:prop:overlap-free:2a}$: 
By \cref{constrsync:eqn:sigmaL-in-suffpref} and the assumption we have $\sigma_{\{v\}}(u) \in \suff{u} \cap \pref{v} \subseteq \{\varepsilon, u, v\}$.
From \cref{constrsync:prop:overlap-free:2} we also get $\suff{u}\cap \pref{u} \subseteq \{\varepsilon, u\}$.

$\ref{constrsync:prop:overlap-free:2a} \Rightarrow \ref{constrsync:prop:overlap-free:2aa}$: 
Suppose $\sigma_{\pref{v}\setminus\{v\}}(u) = w$ for some $w \notin \{\varepsilon, u\}$.
Then we have
\[
\sigma_{\{v\}}(u) = 
\begin{cases}
u & \text{if $u$ is a prefix of~$v$}, \\
w & \text{otherwise.}
\end{cases}
\]
In the first case, we have $w \in \suff{u}\cap \pref{u}$, in the second 
case we have $\sigma_{\{v\}}(u) \notin \{\varepsilon, u, v\}$.

$\ref{constrsync:prop:overlap-free:2aa} \Rightarrow \ref{constrsync:prop:overlap-free:3a}$:
Suppose we have $uv \in L$ for non-empty and proper prefixes $u,v$ of words from~$L$.
Then there exists $w \in \Sigma^+$ such that $vw \in L$.
We have $\sigma_{\pref{vw}\setminus\{vw\}}(uv) = xv$ for some $x \in \Sigma^*$.
By assumption, $x = u$. By \cref{constr_sync:eqn_sigma_pref1}, in this case, $uv$ is a prefix of $vw$.
Thus, $uv$ has $v$ as a proper prefix and 
we find that $\sigma_{\pref{uv}\setminus\{uv\}}(uv)$
has $v$ as a suffix.
Since $uv$ is not a prefix of a word from $\pref{uv}\setminus\{uv\}$,
we also have $\sigma_{\pref{uv}\setminus\{uv\}}(uv) \ne uv$.
Thus \cref{constrsync:prop:overlap-free:2aa} does not hold.

$\ref{constrsync:prop:overlap-free:3a} \Rightarrow \ref{constrsync:prop:overlap-free:1}$: 
 If~$L$ is not overlap-free, then $uv, vw \in L$ for some non-empty $u,v,w \in \Sigma^+$ where
 $u,v$ are non-empty and proper prefixes of words from~$L$.
 Thus, \cref{constrsync:prop:overlap-free:3a} does not hold.
\end{proof}

\begin{remark}
	Note that $\sigma_{\pref{v}\setminus\{v\}}(u) = \varepsilon$
	implies $\sigma_{\{v\}}(u) \in \{\varepsilon, v\}$,
	but the converse implication does not hold. For example, 
	we have $\sigma_{\{aa\}}(baa) = aa$, but $\sigma_{\{a\}}(baa) = a$.
\end{remark} 

\begin{corollary} \label{constrsync:rem:overlap-free}
	Let $u \in \Sigma^*$. Then the following are equivalent.
	\begin{enumerate}
	\item $\{u\}$ is overlap-free.
	
	\item $\sigma_{\pref{u}\setminus\{u\}}(u) = \varepsilon$.
	
	\item $\suff{u}\cap \pref{u} \subseteq \{\varepsilon, u\}$.
	\end{enumerate}
\end{corollary}

\begin{proposition} \label{constr_sync:cor:sigma_L_prefix-eps}
	Let $L \subseteq \Sigma^+$ and $u \in L$ where $L$ contains at least two words.
	Then we have $\sigma_{\pref{L} \setminus \{u\}}(u) = \varepsilon$
	if and only if $\{u\}$ is overlap-free and $\sigma_{L \setminus \{u\}}(u) = \varepsilon$.
\end{proposition}
\begin{proof}
	First, suppose $\sigma_{\pref{L} \setminus \{u\}}(u) = \varepsilon$.
	Since $u \in L$ we have $\pref{L}\setminus \{u\} = ( L \setminus \{u\} ) \cup ( \pref{u} \setminus \{u\})$.
	Thus, by \crefitem{constr_sync:prop:sigma_L_prefix}{constr_sync:prop:sigma_L_prefix:7},
	we find $\sigma_{L \setminus\{u\}}(u) = \varepsilon$
	and $\sigma_{\pref{u}\setminus\{u\}}(u) = \varepsilon$, and thus~$u$ is overlap-free by \cref{constrsync:rem:overlap-free}.
	Conversely, if $\{u\}$ is overlap-free and $\sigma_{L\setminus\{u\}}(u) = \varepsilon$,
	then by \cref{constrsync:rem:overlap-free,constrsync:eqn:sigma_prefL} 
	we obtain $\sigma_{\pref{L} \setminus \{u\}}(u) = \varepsilon$.
\end{proof}

\begin{example}
	Let $L = \{a, ab,ba\}$. Then $L$ is not overlap-free. 
	Note that here we have $uv \notin L$ for all $u,v \in \pref{L} \setminus L = \{\varepsilon, b\}$.
\end{example}

Suppose $uv, vw \in L$ for non-empty $u,v,w \in \Sigma^+$.
Then $uvw \in \left(\bigcup_{x \in L} \pref{x}\setminus\{x\} \right) L \cap L\Sigma^+$. Thus, if the latter
set is empty, then~$L$ is overlap-free. The following statement also shows that when~$L$
is a prefix, this condition characterizes overlap-freeness of~$L$.

\begin{proposition} \label{constrsync:prop:overlap-free-prefix}
Let $L \subseteq \Sigma^+$ be non-empty. Then the following are equivalent.
\begin{enumerate}		
\item \label{constrsync:prop:overlap-free-prefix:1xx} $L$ is an overlap-free prefix code.

\item \label{constrsync:prop:overlap-free-prefix:2xx} For all $u,v \in L$ we have $\suff{u} \cap \pref{v} \subseteq \{\varepsilon, v\}$.

\item  \label{constrsync:prop:overlap-free-prefix:2ax} For all $u,v \in L$ we have $\sigma_{\pref{u}\setminus\{u\}}(v) = \varepsilon$.		

\item  \label{constrsync:prop:overlap-free-prefix:2bx} For all $u \in L$, the language $\{u\}$ is overlap-free and 
$\sigma_{\{v\}}(u) \in \{\varepsilon, v\}$ for distinct $u,v \in L$.

\item \label{constrsync:prop:overlap-free-prefix:8}  For all $u \in L$, $\{u\}$ is overlap-free
and $\sigma_{L\setminus\{u\}}(u) \in (L \setminus\{u\}) \cup \{\varepsilon\}$
if $L \setminus\{u\} \ne \emptyset$.		

\item \label{constrsync:prop:overlap-free-prefix:5ax} For all $u \in L$ we have 
$\sigma_{\pref{L} \setminus \{u\}}(u) \in ( L \setminus \{u\} ) \cup \{\varepsilon\}$.

\item \label{constrsync:prop:overlap-free-prefix:3yx} 
We have $uv \notin L$ for all $u,v \in \Sigma^+$ such that $u \in L$ or $v \in \pref{L} \setminus L$.

\item \label{constrsync:prop:overlap-free-prefix:2axy} $L$ is a prefix code
such that $\left(\bigcup_{u \in L} \pref{u}\setminus\{u\} \right) L \cap L\Sigma^+ = \emptyset$.
\end{enumerate}
\end{proposition}
\begin{proof}	
$\ref{constrsync:prop:overlap-free-prefix:1xx} \Rightarrow \ref{constrsync:prop:overlap-free-prefix:2xx}$: 
This follows from \cref{constrsync:prop:overlap-free} and the fact that 
$u \notin \pref{v}$ for distinct $u,v \in L$ when~$L$ is a prefix code.
	
$\ref{constrsync:prop:overlap-free-prefix:2xx} \Rightarrow \ref{constrsync:prop:overlap-free-prefix:2ax}$: 
We have $\sigma_{\pref{u}\setminus\{u\}}(v) \in \suff{v} \cap ( \pref{u} \setminus\{u\} ) \subseteq \{\varepsilon\}$ by assumption and the definition of~$\sigma_L$.
	
$\ref{constrsync:prop:overlap-free-prefix:2ax} \Rightarrow \ref{constrsync:prop:overlap-free-prefix:2bx}$:
By \cref{constrsync:rem:overlap-free}, we find that $\{u\}$ is overlap-free for each $u \in L$.
Next, for $u,v \in \Sigma^*$ we have
\[
\sigma_{\{v\}}(u) = \begin{cases}
	\sigma_{\pref{v}\setminus\{v\}}(u) & \text{if } v \notin \suff{u}, \\
	v & \text{otherwise.}	
\end{cases}
\]
Thus, we get $\sigma_{\{v\}}(u) \in \{\varepsilon,v\}$
if $u,v \in L$ by \cref{constrsync:prop:overlap-free-prefix:2ax}.

$\ref{constrsync:prop:overlap-free-prefix:2bx}\Rightarrow\ref{constrsync:prop:overlap-free-prefix:8}$: We have $\sigma_{L \setminus \{u\}}(u) = \sigma_{\{v\}}(u)$
for some $v \in L$ with $v \ne u$. By assumption, $\sigma_{\{v\}}(u) \in \{\varepsilon,v\}$
and this gives the claim.

$\ref{constrsync:prop:overlap-free-prefix:8}\Rightarrow\ref{constrsync:prop:overlap-free-prefix:5ax}$: This follows from \cref{constrsync:eqn:sigma_prefL,constrsync:rem:overlap-free}.

$\ref{constrsync:prop:overlap-free-prefix:5ax} \Rightarrow \ref{constrsync:prop:overlap-free-prefix:3yx}$:
Let $u$ be a shortest non-empty word such that $uv \in L$
for some $v \in \Sigma^+$
with $u \in L$ or $v \in \pref{L}\setminus L$.
If $u \in L$,
then $\sigma_{\pref{L}\setminus\{u\}}(u) = u$.
Next, suppose $v \in \pref{L}\setminus L$.
Then $\sigma_{\pref{L}\setminus\{uv\}}(uv) = xv$
for some $x \in \Sigma^*$.
Assume $\sigma_{\pref{L}\setminus\{uv\}}(uv) \in L$.
Since $v \notin L$, this gives $x \ne \varepsilon$.
By the minimal choice of~$u$, we find $x = u$.
Thus we have $\sigma_{\pref{L}\setminus\{uv\}}(uv) = uv$.
This gives the claim.

$\ref{constrsync:prop:overlap-free-prefix:3yx} \Rightarrow \ref{constrsync:prop:overlap-free-prefix:2axy}$:
If $L$ is not a prefix code, then $uv \in L$
for some $u \in L$ and $v \in \Sigma^*$.
Next, suppose~$L$ is a prefix code
and $uv = wx$ for $u$ being a proper prefix of some word from~$L$
and where $v,w \in L$ and $x \in \Sigma^+$.
Since~$L$ is a prefix code,
we have $w = uy$ and $v = yx$ for some $y \in \Sigma^+$
where $y \notin L$. Then we have $uy \in L$ with $y \in \pref{L}\setminus L$.

$\ref{constrsync:prop:overlap-free-prefix:2axy} \Rightarrow 
\ref{constrsync:prop:overlap-free-prefix:1xx}$:
Suppose $uv \in L$ and $vw \in L$ for some $u,v,w \in \Sigma^+$.
Since~$L$ is a prefix code, we have $u \notin L$.
Thus 
$uvw \in \left(\bigcup_{u \in L} \pref{u}\setminus\{u\} \right) L \cap L\Sigma^+$.
\end{proof}

\begin{remark}
	Let $L = \{ aabb, bba, abb \}$ and $u = aabb$.
	Then we have $\sigma_{\pref{L} \setminus \{u\}}(u) = abb$
	where $abb \in L \setminus \{u\}$,
	but $\sigma_{\{bba\}}(u) = bb \notin \{\varepsilon, bba\}$.	
	Thus the equivalences of \cref{constrsync:prop:overlap-free-prefix:2bx} 
	and \cref{constrsync:prop:overlap-free-prefix:5ax}
	in \cref{constrsync:prop:overlap-free-prefix}
	do not hold for a fixed $u \in L$, but only when quantified
	over all~$u \in L$.	
\end{remark}

\begin{corollary}\label{lem:overlap-free-prefix}
	Let $C \subseteq \Sigma^+$ be a prefix code.
	Then $C$ is overlap-free if and only if $(\pref{C} \setminus C) C \cap C \Sigma^+ = \emptyset$.
\end{corollary} 
\begin{proof}
	Let~$C$ be a prefix code.
	Then its set of proper prefixes
	$\bigcup_{u \in L} \pref{u}\setminus\{u\}$
	equals $\pref{C}\setminus C$.
	Thus the claim follows from \cref{constrsync:prop:overlap-free-prefix}.
\end{proof}

\begin{remark}
	Let $L \subseteq \Sigma^*$ be non-empty and $u \in \Sigma^*$.
	We have
	\[
	\sigma_{\pref{L}\setminus \{u\}}(u) \notin \bigcup_{v \in L} ( \pref{v} \setminus \{\varepsilon,v\} )
	\Rightarrow
	\sigma_{\pref{L} \setminus \{u\}}(u) \in ( L \setminus \{u\} ) \cup \{\varepsilon\}.
	\]
	The converse implication holds if no proper suffix of~$u$ contained in~$L$
	is a proper prefix of some other word from~$L$ (which holds if~$L$ is a prefix code).
	This yields that \cref{constrsync:prop:overlap-free-prefix:5ax} 
	from \cref{constrsync:prop:overlap-free-prefix}
	can also be stated by saying the image of the Levenshtein-Romanov mapping
	is not in $\bigcup_{v \in L} ( \pref{v} \setminus \{\varepsilon,v\} )$
	since either condition, when quantified over all words $u \in L$,
	yields that~$L$ is a prefix code.
\end{remark}

\begin{proposition} \label{constrsync:prop:overlap-suffix}
Let $L \subseteq \Sigma^+$. Then the following are equivalent.
\begin{enumerate}
\item \label{constrsync:prop:overlap-suffix:1} $L$ is an overlap-free suffix code.

\item \label{constrsync:prop:overlap-suffix:2} For all $u,v \in L$ we have $\suff{u} \cap \pref{v} \subseteq \{\varepsilon, u\}$.

\item \label{constrsync:prop:overlap-suffix:3} $\{u\}$ is overlap-free for $u \in L$ and $\sigma_{\{v\}}(u) \in \{\varepsilon, u\}$ for $u,v \in L$.

\item \label{constrsync:prop:overlap-suffix:4} $uv \notin L$ for all $u,v \in \pref{L} \setminus \{\varepsilon\}$.	
\end{enumerate}
\end{proposition}
\begin{proof}
$\ref{constrsync:prop:overlap-suffix:1} \Leftrightarrow \ref{constrsync:prop:overlap-suffix:2}$: This follows from \cref{{constrsync:prop:overlap-free:2}} of \cref{constrsync:prop:overlap-free}
and the fact that for distinct $u,v \in L$ we have $v \notin \suff{u}$ if and only if~$L$ is a suffix code.

$\ref{constrsync:prop:overlap-suffix:1} \Leftrightarrow \ref{constrsync:prop:overlap-suffix:3}$:
This follows from \cref{{constrsync:prop:overlap-free:2a}} of \cref{constrsync:prop:overlap-free}
and \cref{constr_sync:eqn_sigma_suff0}.

$\ref{constrsync:prop:overlap-suffix:1} \Leftrightarrow \ref{constrsync:prop:overlap-suffix:4}$: This follows from \cref{constrsync:prop:overlap-free:3a} of
\cref{constrsync:prop:overlap-free}
and the definition of a suffix code.	
\end{proof}

\begin{proposition} \label{constrsync:prop:overlap-free-bifix}
	Let $L \subseteq \Sigma^+$ be non-empty. Then the following are equivalent.	
	\begin{enumerate}
		
		\item \label{constrsync:prop:overlap-free-bifix:1} $L$ is an overlap-free bifix code.
		
		\item \label{constrsync:prop:overlap-free-bifix:2} For $u,v \in L$ we have
		\[
		\suff{u} \cap \pref{v} = \begin{cases}
			\{\varepsilon\} & \text{if } u \ne v, \\
			\{\varepsilon, u\} & \text{if } u = v.	
		\end{cases}
		\]
		
		\item \label{constrsync:prop:overlap-free-bifix:3} For $u \in L$ we have $\sigma_{\pref{L}\setminus\{u\}}(u) = \varepsilon$.\footnote{As stated
			in \cref{constr_sync:cor:sigma_L_prefix-eps},
			this is equivalent to $\{u\}$ being overlap-free
			and $\sigma_{L\setminus\{u\}}(u) = \varepsilon$
			if $L \setminus \{u\} \ne \emptyset$.}

		\item \label{constrsync:prop:overlap-free-bifix:4} 
		We have $uv \notin L$ for all $u,v \in \Sigma^+$ with $u \in L$ or $v \in \pref{L}$.
	\end{enumerate}
\end{proposition}
\begin{proof}
$\ref{constrsync:prop:overlap-free-bifix:1} \Leftrightarrow \ref{constrsync:prop:overlap-free-bifix:2}$: This follows from \cref{constrsync:prop:overlap-free,constrsync:prop:bifix-code}.

$\ref{constrsync:prop:overlap-free-bifix:1} \Leftrightarrow \ref{constrsync:prop:overlap-free-bifix:3}$:
First, suppose $L$ is an overlap-free bifix code.
By \crefitem{constrsync:prop:overlap-free-prefix}{constrsync:prop:overlap-free-prefix:5ax}
we have that $\sigma_{\pref{L}\setminus\{u\}}(u)$ is in $( L \setminus \{u\} ) \cup \{\varepsilon\}$.
Since $\sigma_{\pref{L}\setminus\{u\}}(u)$ is a suffix of $u$, 
it cannot be contained in $L \setminus \{u\}$ since~$L$ is a suffix code.
Conversely, by \crefitem{constrsync:prop:overlap-free-prefix}{constrsync:prop:overlap-free-prefix:5ax}
this condition implies that~$L$ is an overlap-free prefix code.
Next, if $uv \in L$ for $u \in \Sigma^+$ and $v \in L$,
then $\sigma_{\pref{L}\setminus\{uv\}}(uv)$ has the non-empty word $v$ as a suffix.
Thus, $L$ must be a suffix code.

$\ref{constrsync:prop:overlap-free-bifix:1} \Leftrightarrow \ref{constrsync:prop:overlap-free-bifix:4}$: 
This follows from \crefitem{constrsync:prop:overlap-free-prefix}{constrsync:prop:overlap-free-prefix:3yx}
and the definition of a suffix code.
\end{proof}

Note that $L \subseteq \Sigma^+$ is a prefix (resp.~suffix, bifix, infix, overlap-free) code/language
if and only if $\{u,v\}$ is a prefix (resp.~suffix, bifix, infix, overlap-free) code/language
for all $u,v \in L$. So, the conditions from the above statements characterizing those language classes
can be written as conditions that must be true for all pairs (of possibly equal) words in $L$.

The following further characterizations were given in~\cite{DBLP:journals/tcs/JurgensenS21}
for non-empty $L \subseteq \Sigma^+$: 
\begin{enumerate}
	\item $\pref{L}\cap \Sigma^*L = L$ if and only if $\sigma_L(u) \in L$ for all $u \in \Sigma^*L$,
	
	\item $L$ is an overlap-free prefix code if and only if $\sigma_L(uv) = \sigma_L(v)$
	for all $u \in L$ and $v \in \Sigma^+$,
	
	\item $L$ is a suffix code such that $\pref{L} \cap \Sigma^*L = L$ if and only if $\sigma_L(uv) = v$
	for all $u \in \Sigma^+$ and $v \in L$ (this condition also appears in \crefitem{constrsync:prop:infix-code}{constrsync:prop:infix-code:6}).
\end{enumerate}

In the next section we will investigate overlap-free infix codes, i.e., solid codes, in more detail.

We always have $L \cup \{\varepsilon\} \subseteq \pref{L} \cap \suff{L}$
for any language $L \subseteq \Sigma^*$. Lastly, we connect
the converse inclusion to overlap-freeness.

\begin{proposition}  \label{constrsync:prop:pref-suff-overlap-free-inclusion-eqn}
	Let $C \subseteq \Sigma^+$ be a prefix or a suffix code.
	If $\pref{C} \cap \suff{C} \subseteq C \cup \{\varepsilon\}$,
	then $C$ is overlap-free.
\end{proposition}
\begin{proof}
	Let $uv, vw \in C$ for non-empty $u,v,w$.
	Then we have $v \in \pref{C} \cap \suff{C}$
	and since~$C$ is a prefix or suffix code, we have $v \notin C$, i.e.,
	$\pref{C} \cap \suff{C}$
	is not a subset of~$C$.
\end{proof}

\begin{proposition} \label{constrsync:prop:overlap-free-inclusion-eqn}
	Let $L \subseteq \Sigma^*$.
	If $L$ is overlap-free,
	then $\pref{L} \cap \suff{L} \subseteq L \cup \{\varepsilon\}$.
\end{proposition}
\begin{proof}
	Let $u \in \pref{L}\cap \suff{L}$ with $u \ne \varepsilon$.
	Then $uv, wu \in L$ for some $v,w \in \Sigma^*$.
	If~$L$ is overlap-free, then $v = \varepsilon$
	or $w = \varepsilon$ and each case gives $u \in L$,
	which shows the claim.	
	%Let $v \in \pref{L}\cap \suff{L}$ with $v \ne \varepsilon$ and $v \notin L$.
	%Then we have $uv, vw \in L$ for some non-empty $u,w \in \Sigma^+$, i.e., $L$ 
	%is not overlap-free.
\end{proof}

Thus, we find that a prefix or suffix code~$C$
is overlap-free if and only if $\pref{C} \cap \suff{C} \subseteq C \cup \{\varepsilon\}$.

\subsection{Characterizing Solid Codes}
\label{constrsync:sec:codes:solid}

%!TEX root = ./../../springer-mono-format/book.tex
%!TEX root = ./../../thesis-trier.tex
 
\IN{strongly regular codes}
Solid codes provide outstanding fault-tolerance when used for information transmission through a noisy channel w.r.t.~symbol substitutions and synchronization errors. In this context,
they were introduced by Levenshtein~\cite{Lev64} under the name \emph{strongly regular codes}, 
and further investigated by Romanov~\cite{Rom66}.
Shyr and Yu~\cite{ShyrYu90} introduced solid codes to investigate disjunctive domains, 
but did not recognize the connection to the earlier work by Levenshtein and Romanov. 
Later, Jürgensen and Yu~\cite{DBLP:journals/eik/JurgensenY90} observed that strongly 
regular codes describe precisely the solid codes.
For further information, we refer to the surveys~\cite{DBLP:journals/iandc/Jurgensen08,DBLP:reference/hfl/JurgensenK97},
or the introductory section of~\cite{DBLP:journals/tcs/JurgensenS21}. 
Recall that \cref{constrsync:ex:binarysolidcodeitem,constrsync:ex:constr-solid-code}
from \cref{constrsync:ex:binarysolidcode} gave an infinite binary solid code 
and a method to construct a solid codes, respectively.

In our investigation on the constrained synchronization problem and synchronization conditions
for automata obtained in reductions, we discovered (to the best of our knowledge)
new characterizations of solid codes, which we present in the following
(see the discussion section,
\cref{constr_sync:sec:gen_results:discussion}, page~\pageref{discussion:strongly-self-sync-codes}, 
for some background on how this characterization arose in the context
of the constrained synchronization problem).

Here, first, in \cref{thm:char-solid-codes}, we recap known characterizations. 
\Cref{thm:solidcode} collects further characterizations of solid codes
which are new and connect solid codes to various language equations.
Lastly, \cref{thm:solidcode-char2,cor:solidcode-char3,constrsync:prop:sigma-solid-code-char} 
give two other new characterizations.

\begin{theorem} \label{thm:char-solid-codes}
Let $L \subseteq \Sigma^+$ with $L \ne \emptyset$. Then the following are equivalent.
\begin{enumerate}
	
\item  \label{thm:char-solid-codes:1} $L$ is a solid code.
	
\item  \label{thm:char-solid-codes:2} (Jürgensen \& Yu~\cite{DBLP:journals/eik/JurgensenY90}) For every word $w \in \Sigma^*$ 
 we have $w = u_0 w_1 u_1 w_2 u_2 \cdots w_n u_n$
 for some $n \ge 0$ and unique words $u_0, u_1, \ldots, u_n \in \Sigma^* \setminus \Sigma^* L \Sigma^*$
 and $w_1, w_2, \ldots, w_n \in L$.
	
\item \label{thm:char-solid-codes:3} (Staiger~\cite{DBLP:journals/tcs/JurgensenS21}, Romanov~\cite{Rom66}, Levenshtein~\cite{Lev64})
 There exists a state-invariant decoder without look-ahead for $L$ (see \cite[Def.~9]{DBLP:journals/tcs/JurgensenS21}
 for a definition of these decoders for finite solid codes and extensions to infinite solid codes).

\item \label{thm:char-solid-codes:4} (Liu, Guo \& Shum~\cite[Thm.~2.4]{Liu2021}) $L$ is an infix code and $\{ u_1\cdots u_n \}$ is overlap-free for 
 all sequences of distinct words $u_1, \ldots, u_n \in L$ with $n \ge 1$.\footnote{In~\cite[Thm.~2.4]{Liu2021},
 	if~$L$ is finite, then overlap-freeness is only stated to hold for $\{u_1\cdots u_n\}$
 	if $2 \le n \le |L|$. However, this seems to be a typo,
 	as for example $L = \{ aba, bcb \}$
 	is an infix code
 	and $\{ ababcb\}$ and $\{ bcbaba \}$
 	are each overlap-free, while~$L$
 	is not overlap-free and so not a solid code.} 
 
\item \label{thm:char-solid-codes:5} (Liu, Guo \& Shum~\cite[Thm.~2.4]{Liu2021}) $L$ is an infix code such that $\{u\}$ is overlap-free for all $u \in L$ and $\{uv\}$ is overlap-free for $u,v \in L$ with $u \ne v$.
\end{enumerate}	
\end{theorem}

\begin{remark}
	Since for distinct $u_1, \ldots, u_n \in L$
	we have $u_i \cdots u_n \ne u_1 \cdots u_{n - i + 1}$
	for all $i \in \{1,\ldots,n\}$.
	Thus~$L$ is overlap-free
	if $\{ u_1 \cdots u_n \}$
	is overlap-free for distinct $u_1, \ldots, u_n \in L$.
	Next, if for distinct $u,v \in L$
	we have $u = xy$ and $v = yz$
	for non-empty $x,y,z \in \Sigma^+$,
	then $vu = yzxy$ and so $\{vu\}$
	is not overlap-free.
	Thus, if $\{u\}$ for $u \in L$ and $\{uv\}$
	for distinct $u,v \in L$ are overlap-free,
	then~$L$ is overlap-free.
	This shows that the additional conditions in Items~\ref{thm:char-solid-codes:4} and~\ref{thm:char-solid-codes:5} of \cref{thm:char-solid-codes}, beyond being an infix code, are equivalent to overlap-freeness (without any additional assumptions on~$L$).
\end{remark}

\begin{theorem} \label{thm:solidcode}
Let $L \subseteq \Sigma^+$. Then the following are equivalent:
\begin{enumerate}
\item  \label{thm:solidcode:1} $L$ is a solid code.

\item  \label{thm:solidcode:1a} $L$ is a bifix code and $(\pref{L} \setminus L) L \cap \Sigma^* L \Sigma^+ = \emptyset$.

\item  \label{thm:solidcode:1b}  $L$ is a suffix code and $\left( \bigcup_{u \in L} ( \pref{u}\setminus\{u\} ) \right)L \cap \Sigma^*L\Sigma^+ = \emptyset$.

\item \label{thm:solidcode:2} $L$ is an infix code and $(\pref{L} \setminus L) L \cap L \Sigma^+ = \emptyset$.

\item \label{thm:solidcode:3} $L$ is an infix code and $\pref{L} \cap \suff{L} = L \cup \{\varepsilon\}$.	
\end{enumerate}
\end{theorem}
\begin{proof}
$\ref{thm:solidcode:1}\Rightarrow \ref{thm:solidcode:1a}$: A solid code is an infix code, and so a bifix code, and overlap-free.
Suppose $vw = xuy$
with $x \in \Sigma^*$, $y \in \Sigma^+$, $u, w \in L$, and
$v \in \pref{L} \setminus L$.
By \cref{lem:overlap-free-prefix}
we have $|x| > 0$.
For $|x| \ge |v|$ (resp.~$|y| \ge |w|$), we have $u \in \factor{w}$ (resp.~$u \in \factor{v}$).
For $0 < |x| < |v|$ and $0 < |y| < |w|$,
we have $u = v'w'$ for a non-empty
and proper suffix $v'$ (resp.~prefix~$w'$)
of $v$ (resp.~$w$), so that $L$ is not overlap-free.
Hence, for $(\pref{L} \setminus L) L \cap \Sigma^+ L \Sigma^+  \ne \emptyset$,
we get that $L$ is not a solid code, which gives the claim.

$\ref{thm:solidcode:1a}\Rightarrow \ref{thm:solidcode:1b}$: Since~$L$ is a prefix code,
 its set of proper prefixes equals $\pref{L}\setminus L$, which gives the claim.

$\ref{thm:solidcode:1b}\Rightarrow \ref{thm:solidcode:2}$: Suppose~$L$ is a suffix code as stated.
Assume we have $uvw \in L$ with $u \in \Sigma^*$, $v \in L$,
and $w \in \Sigma^+$.
Let $x$ be a proper prefix of some word in~$L$.
Then $xuvw \in \Sigma^* L \Sigma^+$
and this is excluded by assumption.
Thus, $L$ is an infix code.
The assumption yields that $( \pref{L}\setminus L ) L \cap L \Sigma^+ = \emptyset$.

$\ref{thm:solidcode:2} \Rightarrow \ref{thm:solidcode:3}$: 	By \cref{lem:overlap-free-prefix}, $L$ is overlap-free.
Thus, \cref{constrsync:prop:overlap-free-inclusion-eqn}
gives $\pref{L} \cap \suff{L} \subseteq L \cup \{\varepsilon\}$.
The inclusion
$L \cup \{\varepsilon\} \subseteq \pref{L} \cap \suff{L}$
holds for any language.	

$\ref{thm:solidcode:3} \Rightarrow \ref{thm:solidcode:1}$: 
By \cref{constrsync:prop:pref-suff-overlap-free-inclusion-eqn}, $L$
is overlap-free and thus a solid code.	
\end{proof}

\begin{remark}
	As far as we are aware of,
	the condition $(\pref{L} \setminus L) L \cap \Sigma^* L \Sigma^+ = \emptyset$
	was never investigated in the literature before (apart
	from our own work~\cite{DBLP:conf/fct/Hoffmann21} 
	whose results are presented and extended
	in the present thesis).
	Note that $\pref{L}L \subseteq L\Sigma^+$ holds for every language $L \subseteq \Sigma^*$.
\end{remark}

We mention an implication for languages  $L \subseteq \Sigma^*$
such that $(\pref{L} \setminus L) L \cap \Sigma^* L \Sigma^+ = \emptyset$.

\begin{proposition}
	Let $L \subseteq \Sigma^*$
	such that $(\pref{L} \setminus L) L \cap \Sigma^* L \Sigma^+ = \emptyset$.
	Then $\pref{L}\subseteq L$
	or $uvw \in L$ with $u,v \in \Sigma^*$ and $v \in L$ implies $w = \varepsilon$.
\end{proposition}
\begin{proof}
	Suppose we have  $x \in \pref{L} \setminus L$
	for some $x \in \Sigma^*$.
	Let $v \in L$
	with $uvw \in L$ for some $u,w \in \Sigma^*$.
	Assume $|w| > 0$.
	Then we have $xuvw \in (\pref{L}\setminus L) L \cap \Sigma^* L \Sigma^+$.
	Thus, we must have $|w| = 0$.
\end{proof}

\begin{theorem} \label{thm:solidcode-char2}
	Let $L \subseteq \Sigma^+$.
	Then $L$ is a solid code
	if and only if for all $u \in \pref{L}$ and $v \in L$
	such that $uv = xwy$ for  $x \in \Sigma^*$ and $y \in \Sigma^+$ and $w \in L$
	we have $x = \varepsilon$ and $u = w$.			
\end{theorem}
\begin{proof}
	First, suppose $L$ is a solid code.
	Let $u \in \pref{L}$, $v, w \in L$,
	and $uv = xwy$ for $x \in \Sigma^*$ and $y \in \Sigma^+$.	
	By \crefitem{thm:solidcode}{thm:solidcode:1a}, we must have $u \in L$.
	By \crefitem{constr_sync:prop:selfsync_implies_avoiding}{constr_sync:prop:item:sc-in-cf}, 
	$L$ is comma-free.
	So $u, v \in L$	imply $x = \varepsilon$ and, as $L$ is a prefix code, $w = u$.	
	Conversely, suppose the stated condition holds.
	Then, clearly, $L$ is comma-free.
	Next, suppose $u \in \pref{L} \setminus L$ and $v,w \in L$
	with $uv = xwy$ for some $x \in \Sigma^*$ and $y \in \Sigma^+$.
	By assumption, $x = \varepsilon$ and $u = w$.
	However, as $u \notin L$ we have $u \ne w$.
	So we must have $(\pref{L} \setminus L)L \cap L \Sigma^+ = \emptyset$.
\end{proof}

\begin{corollary} \label{cor:solidcode-char2}
	Let $C \subseteq \Sigma^+$ be a solid code.
	Then $\pref{C}C \cap \Sigma^+ C \Sigma^+ = \emptyset$.
\end{corollary}

The condition $\pref{C}C \cap \Sigma^+ C \Sigma^+ = \emptyset$
clearly implies that~$C$ is comma-free,
but not every comma-free code fulfills it.
For example, for the comma-free code $C = \{aba,cd,bc\}$
we have $abcd \in (\pref{C}\setminus C)C \cap \Sigma bc \Sigma$.

The next examples show that $\pref{C}C \cap \Sigma^+ C \Sigma^+ = \emptyset$
does not characterize solid codes.

\begin{example}
	\begin{enumerate}
		\item For the comma-free code $C = \{abb,aba\}$ from \crefitem{constrsync:ex:binarysolidcode}{constrsync:ex:comma-free-not-solid}
		that is not a solid code we have $\pref{C}C \cap \Sigma^+ C \Sigma^+ = \emptyset$.
		
		\item Let $C = \{ab,bc\}$. Then $C$ is not a solid code,
		but we have $\pref{C}C \cap \Sigma^+ C \Sigma^+ = \emptyset$.
	\end{enumerate} 	
\end{example}

\begin{proposition} \label{constrsync:prop:prefLintersect}
	Let $L \subseteq \Sigma^+$.
	Then $\pref{L}L \cap \Sigma^+ L \Sigma^+ = \emptyset$
	if and only if for all $u \in \pref{L}$ and $v \in L$
	such that $uv = xwy$ for $x \in \Sigma^+$ and $y \in \Sigma^*$ and $w \in L$
	we have $y = \varepsilon$ and $w = v$.
\end{proposition}
\begin{proof}
	First, suppose $\pref{L}L \cap \Sigma^+ L \Sigma^+ = \emptyset$.
	In this case, we must have $y = \varepsilon$ since
	otherwise $uv \in \Sigma^+ L \Sigma^+$.
	Since~$L$ is a suffix code, this implies $w = v$.
	
	Conversely, the condition 
	implies $\pref{L}L \cap \Sigma^+ L \Sigma^+ = \emptyset$.	
\end{proof}

\begin{corollary} \label{cor:solidcode-char3}
	Let $L \subseteq \Sigma^+$.
	Then $L$ is a solid code
	if and only if for all $u \in \pref{L}$ and $v \in L$,
	if $uv = xwy$ for  $w \in L$ and $x,y \in \Sigma^*$,
	then $x = \varepsilon$ and $w = u$
	or $y = \varepsilon$ and $w = v$.
\end{corollary}
\begin{proof}
	\Cref{thm:solidcode-char2,cor:solidcode-char2,constrsync:prop:prefLintersect} show 
	that a solid code fulfills the condition.
	Conversely, it implies the condition
	from \cref{thm:solidcode-char2} if $|y| > 0$
	and thus that~$L$ is a solid code.
\end{proof}

Lastly, we give a characterization using the Levenshtein-Romanov mapping. It basically combines results
from \cref{constrsync:sec:codes:levensthein-romanov}.

\begin{proposition} \label{constrsync:prop:sigma-solid-code-char}
	Let $L \subseteq \Sigma^+$ be non-empty.
	Then $L$ is a solid code if and only if
	for all $u,v \in \Sigma^*$ the following is true:
	\begin{enumerate}
		\item \label{constrsync:prop:sigma-solid-code-char:1} 
		$\sigma_L(uv) = uv$ implies $u \notin L$ or $v = \varepsilon$,
		
		\item \label{constrsync:prop:sigma-solid-code-char:2} $\sigma_L(uv) \in L$ implies $v \notin \pref{L}\setminus L$ or $v = \varepsilon$,
		
		\item \label{constrsync:prop:sigma-solid-code-char:3} $\sigma_L(uv) \ne v$ implies $v \notin L$. 	
	\end{enumerate}
\end{proposition}
\begin{proof}
	First, suppose~$L$ is a solid code.	
	Assume $\sigma_L(uv) = uv$.
	Since~$L$ is a prefix code,
	by \cref{constrsync:prop:prefix-code:4} of \cref{constrsync:prop:prefix-code}
	we find that~$u \notin L$
	or $v = \varepsilon$.
	Next, let $\sigma_L(uv) \in L$ and $v \in \pref{L}\setminus L$.
	Then $\sigma_L(uv) = wv$ for some $w \in \Sigma^*$
	with $|w| > 0$ since $wv \in L$.
	By overlap-freeness, this yields $v = \varepsilon$.	
	\Cref{constrsync:prop:sigma-solid-code-char:3}
	follows from \cref{constrsync:prop:infix-code:6} of \cref{constrsync:prop:infix-code}
	since~$L$ is an infix code.
		
	Conversely, \cref{constrsync:prop:sigma-solid-code-char:1} 
	yields that~$L$ is a prefix code by \crefitem{constrsync:prop:prefix-code}{constrsync:prop:prefix-code:4}.
	Then, \cref{constrsync:prop:sigma-solid-code-char:3} 
	yields that~$L$ is an infix code by \cref{constrsync:prop:infix-code:6} 
	of \cref{constrsync:prop:infix-code}.
	Next, assume $uv, vw \in L$ for non-empty $u,v,w \in \Sigma^+$.
	We have $v \in \pref{L}\setminus L$ since~$L$ is a prefix code
	and $\sigma_L(uv) = uv$. Thus, \cref{constrsync:prop:sigma-solid-code-char:2} 
	does not hold. So, if all these conditions hold,
	we find that~$L$ is overlap-free.
\end{proof}

\begin{remark}
	By \cref{constrsync:prop:sigma-solid-code-char:3} 
	from \cref{constrsync:prop:sigma-solid-code-char},
	for a solid code~$C$ and $v \in C$,
	we have $\sigma_C(uv) = v$ for $u \in \Sigma^*$.
	By \cref{constrsync:prop:sigma-solid-code-char:2},
	we also have $\sigma_C(uw) \notin L$
	for every proper and non-empty prefix $w$
	of a word from~$C$.
	However, note that we do not have $\sigma_C(uw) = w$
	in this case. For example, for the solid code $C = \{abb, aac\}$
	we have $\sigma_C(uw) = uw$ for $u = a$
	and $w = a$.
\end{remark}